\documentclass[11pt]{article}
\usepackage[margin=1in]{geometry}
\usepackage{amsmath,amssymb}
\usepackage{graphicx}
\usepackage{booktabs}
\usepackage{multirow}
\setlength{\emergencystretch}{3em}
\usepackage{listings}
\usepackage{xcolor}
\usepackage[hidelinks]{hyperref}
\renewcommand\labelitemi{$\bullet$}
\lstset{basicstyle=\ttfamily\scriptsize,breaklines=true,columns=fullflexible,keepspaces=true,frame=single,literate={\\}{{\char92}}1 {\$}{{\char36}}1 {\{}{{\char123}}1 {\}}{{\char125}}1 {−}{{-}}1 {×}{{x}}1 {±}{{+-}}1 {·}{{.}}1 {≤}{{<=}}1 {≥}{{>=}}1 {–}{{--}}1 {—}{{---}}1}
\newcommand{\grade}[1]{\textsf{\small[#1]}}
\title{What the Neutrino State Must Keep:\\ Hidden Angular Structure, the Flux--Shear Link\\ and the Threshold for Primordial Black Hole Formation}
\author{Jeromie Beasley}
\date{14 November 2026 \quad (version 12 of \emph{When the Neutrinos Leave})}
\begin{document}
\maketitle

\noindent\emph{AI-assisted. The ideas and research direction are the author's; AI tools (Claude, Anthropic) were used for literature search, computation, drafting and typesetting. The numbers come from the scripts, run logs and outputs printed in Appendix~B. The simulations are runs of S.~McCormick's public code on a cloud computer, not experiments.}

\paragraph{Physics in one paragraph.} A neutrino gas reaching decoupling can carry a hidden angular record, invisible to its density, flux, stress and to Einstein's equations. Measured in McCormick's general-relativistic code, that record moves the black-hole threshold by up to $2.8\times10^{-3}$ from the unperturbed run, with a sign set by its orientation. In a Gaussian universe, at the planted amplitude, this changes the number of black holes near $6\times10^5\,M_\odot$ by 5--47\%, about three to eleven times what the flux--shear split does in McCormick's physical model, and even a randomly oriented record raises the abundance, by up to 8\% along the Gramian's leading direction. The planted amplitude is engineered, half of each cell's positivity margin; how the shift scales with amplitude is untested. A universe that starts blank carries no such record; a self-recurring universe, which relaxes and lets go again, does. The black-hole mass function across the neutrino step can tell the two apart. The neutrinos themselves sit on the same wall: their heaviest mass step lies at $\Gamma^2=\tfrac12$, Alikhanov's relation read as the wall, within $0.41\sigma$ on the current data, which needs normal ordering and a lightest mass below about 1.5\,meV (Section~\ref{sec:physics}).

\begin{abstract}
McCormick has shown that neutrinos which stream freely out of a collapsing overdensity raise the threshold for primordial black hole formation, by about $0.235\,R_\nu$ in the linear compaction, and that the effect fades when the neutrinos decouple gradually. This paper asks what sets that fading. First, the neutrino collision operator is read exactly as a cut between what is kept and what streams. All of its non-normality, $\tfrac18\lVert[L,L^\ast]\rVert^2=\kappa^2k^2/3$, lies on the single link between density and flux. With separate rates for flux and shear it is $k^2\kappa_1^2/3+\tfrac{4k^2}{15}(\kappa_2-\kappa_1)^2$, and the weight across the cut equals the sound speed, $\sqrt{E_B}=k/\sqrt3$. The kept sector is the charge-diffusion channel of relaxation kinetic theory, whose mode exists only for $\kappa>2k/\pi$ (Romatschke); its non-orthogonality is $4\phi^2/\sin^22\phi$, $\phi=k/\kappa$, which diverges there. Second, McCormick's own code is run on cases outside his paper, with every prediction and kill line written down before the runs. A broad profile that collapses near $\epsilon\approx7$--$9$ keeps the same share of the neutrino effect as a Gaussian that collapses near $\epsilon\approx4$, or less, so the collapse time alone does not set the clock. A steeper decoupling history, $\kappa/H=(\epsilon_{\rm dec}/\epsilon)^6$, gives $G=0.648$, $0.317$ and $0.154$ at $\epsilon_{\rm dec}=1$, $2$ and $3$, against $0.683$, $0.431$ and about $0.30$ for the standard history: how the neutrinos leave matters, not only when. A kept-sector memory form, the population the density keeps under the decoupling operator, fitted once to McCormick's curve and then frozen, predicted those three histories within its kill line of $0.08$ while a history-blind rule failed. Third, the memory has a location. Two neutrino states with the same density, flux and anisotropic stress, differing only in finer angular structure, give later density, flux and stress that differ by several per cent when the difference is set at or after decoupling; set deep in the coupled era it is erased, by the factor $\exp[-\epsilon_{\rm dec}^n/(n\tau_s^n)]$. An exact block elimination puts the whole effect in one hidden-source term, and the kept population obeys $(\dot A)^2\le4E_BA(1-A)$. No instantaneous closure on density, flux and stress alone can carry this memory. Fourth, the anisotropic stress sets the threshold. With the flux relaxing at $\kappa_1$ and the shear and every higher moment at $\kappa_2=r\kappa_1$, eleven thresholds from $r=1/16$ to $16$ show that the surviving share is not set by the connection energy, which is stationary at $r=1$, but by the direction of the flux--shear mismatch: $G$ runs from $0.777$ to $0.129$, and its offset $\ln[(1-G)/G]$ moves close to a straight line in the link rapidity $\tfrac12\ln r$, at a rate of $1.10\pm0.01$ inside the window where the split operator is a Markov generator; with the collision rate evaluated at every Runge--Kutta stage the two inner pairs agree to $0.004$, and the $\pm0.01$ is an error budget, not a converged value or a statistical interval; a third grid level shrinks the resolution shifts of the two thresholds to $0.42$ and $0.49$ of the previous step and gives $1.101$, against $1.102$ at the second level, with a bracket range of $1.094$--$1.109$. Damping the shear alone, with every higher moment left at the flux rate, gives a larger rate, $1.19$, where any rule built from link-weight ratios predicts less than $0.1$: the threshold reads the shear itself. The frozen memory form misses this effect entirely, and so does the intrinsic displacement of the linear kept-sector state at readout, which the two collision models move by almost the same angle: of the rulers tested, the law lives in the operator's link rapidity, and the linear kept state at readout does not hold the difference between the models. Halving the neutrino fraction to $R_\nu=0.2$ leaves $G$ at $0.409$, against $0.425$ at $R_\nu=0.409$ measured with the same code. Fifth, the hidden sector is weighed. Its readout Gramian has no kernel, so the predictive quotient is the whole hidden space and the Gramian's spectrum ranks what a reduced state must keep: $10^{-7}$ for structure planted in the coupled era, $0.37$ at decoupling, always led by $\ell=3$. Two collision histories with equal optical depth and equal integrated connection energy, the rate reversed in time, differ in the memory form by up to $0.16$: the order of the history is part of the record. Sixth, the hidden structure moves the threshold itself. In McCormick's full general-relativistic collapse, two neutrino states identical in density, flux, stress, lapse, constraints and every fluid and metric array, differing only in hidden angular structure planted at decoupling, have black-hole thresholds $3.7\times10^{-3}$ apart for $\ell=3$, $2.1\times10^{-3}$ with the opposite sign for $\ell=4$, and $5.5\times10^{-3}$ along the Gramian's leading direction, two hundred bracket widths or more; planted in the coupled era they are $2\times10^{-5}$ apart. The threshold reads angular information that no closure on density, flux and stress holds. Seventh, McCormick's multigroup collision term, with his archived per-group target, is exactly the split link, with $r=1+\kappa_n/\kappa_e$: neutrino self-scattering is the flux--shear split. Removing the split while keeping the isotropy history moves his threshold by $+0.0003$, a third of what the frozen link law predicts; most of what self-scattering does is timing. The algebraic identities, the counterexample and every verdict's arithmetic are machine-checked in Lean (Appendix~C). Two further tests are stated for the e$^\pm$ annihilation epoch and for the mean of the rates.
\end{abstract}

\section{The question}

Primordial black holes heavier than about $10^6\,M_\odot$ form while the neutrinos, which carry 41--49\% of the radiation, stop scattering and stream freely. McCormick evolved these neutrinos with their angle-resolved transport equation inside spherical general relativity and found that free streaming makes collapse harder~\cite{McC}. The threshold on the linear compaction $C_l$, which sets abundance, rises by $\Delta C_{l,c}/C_{l,c}\simeq0.235\,R_\nu$, and the threshold $\delta_c$ on the compaction peak rises by $\simeq0.15\,R_\nu$. For Gaussian fluctuations a perfect-fluid collapse probability $\beta$ becomes about $\beta^{1.18}$ above $6\times10^5\,M_\odot$.

He also ran gradual decoupling, with collision rate $\kappa=H(\epsilon_{\rm dec}/\epsilon)^3$, where $\epsilon=1/(aHr_m)$ is the Hubble radius in units of the perturbation size. Write $G$ for the share of the instantaneous free-streaming shift that survives:
\begin{equation}
G=\frac{\delta_c-\delta_c^{\rm fluid}}{\delta_c^{\rm inst}-\delta_c^{\rm fluid}}.
\end{equation}
His Table~II gives $G=$ 0.959, 0.867, 0.683, 0.431, 0.190 and 0.063 at $\epsilon_{\rm dec}=$ 0.25, 0.5, 1, 2, 4 and 8. He notes that linear theory gives estimates of opposite sign and ``does not determine which diagnostic, if any, predicts the nonlinear threshold''~\cite[Sec.~IV\,C]{McC}.

Two questions follow and are answered here. What in the neutrino transport operator carries the effect (Section~\ref{sec:cut})? And which clock, and which part of the decoupling history, sets $G$ (Sections~\ref{sec:runs}--\ref{sec:memory})? Version~6 adds three more. Which hidden directions does the future read, and must a reduced state keep them (Section~\ref{sec:gramian})? Does hidden angular structure move the threshold itself, in the full collapse (Section~\ref{sec:door})? And what is McCormick's physical collision model in the language of the flux--shear link (Section~\ref{sec:mcc})?

\section{The physics: what the neutrino record does to black holes}\label{sec:physics}

\noindent\emph{This section states what the results of the later sections change in the universe. It adds no new collapse runs. Its estimates use the thresholds measured in Sections~\ref{sec:link}--\ref{sec:mcc} and one standard abundance formula, and each carries its grade.}

\paragraph{The picture.} A neutrino gas carries two kinds of information. The \emph{visible} part is its density, flux and stress: the three numbers a fluid knows, and the only ones Einstein's equations see. The \emph{hidden} part is the finer angular structure, the moments $\ell\ge3$. While the neutrinos collide, the hidden part is wiped. Collisions keep only the isotropic part, and the door of Section~\ref{sec:witness} is shut. When the collision rate falls through the wall $\kappa=2k/\pi$ (Section~\ref{sec:cut}), the kept mode merges into the streaming continuum and the gas lets go. Light-speed streaming is the fastest channel. It leaves first and leaves the lightest trail. Everything slower then follows the trail it left: the hidden structure streams back into flux, stress and density, and the collapse reads it. The threshold for a black hole therefore depends on the angular record the neutrinos carried through decoupling, not only on how much neutrino there is \grade{reading, built on the measured door of Section~\ref{sec:door}}.

\paragraph{Why the sign of the record matters.} Streaming climbs the moment chain one rung at a time (Section~\ref{sec:cut}). Odd hidden structure ($\ell=3$) returns to the visible sector first as \emph{flux}, two rungs down. Depending on its orientation, it returns as a net neutrino push out of the overdensity or into it. Outward flux carries energy away and makes collapse harder; inward flux makes it easier. Even hidden structure ($\ell=4$) returns first as \emph{stress}. Its orientation sets whether the radial pressure exceeds the tangential pressure or falls below it, which resists the infall or eases it. The measurements agree with this ordering:
\begin{itemize}
\item both responses are odd in the amplitude ($\ell=3$: $+1.83$ and $-1.91\times10^{-3}$; $\ell=4$: $-1.08$ and $+1.02\times10^{-3}$, from the unperturbed restart);
\item the odd, flux-first structure moves the threshold more than the even, stress-first structure, by a factor of $1.8$;
\item the Gramian's leading direction, mostly $\ell=3$, moves it most ($+2.69$ and $-2.82\times10^{-3}$).
\end{itemize}
\grade{measured shifts; mechanism a reading, not yet isolated by a run}.

\paragraph{What it does to the number of black holes.} For Gaussian fluctuations the collapse fraction is $\beta=\mathrm{erfc}[\delta_c/(\sqrt2\,\sigma)]$. Tune $\sigma$ so the unperturbed $\delta_c=0.51075$ gives $\beta_0=10^{-6}$, $10^{-10}$ or $10^{-16}$, then insert the measured shifts:
\begin{center}\small\begin{tabular}{lcccc}\toprule
hidden record at decoupling & $\Delta\delta_c$ & $\beta/\beta_0$ ($10^{-6}$) & ($10^{-10}$) & ($10^{-16}$)\\\midrule
$\ell=3$, one orientation & $+1.83\times10^{-3}$ & $0.92$ & $0.86$ & $0.78$\\
$\ell=3$, the other & $-1.91\times10^{-3}$ & $1.10$ & $1.17$ & $1.30$\\
$\ell=4$, one orientation & $-1.08\times10^{-3}$ & $1.05$ & $1.09$ & $1.16$\\
$\ell=4$, the other & $+1.02\times10^{-3}$ & $0.95$ & $0.92$ & $0.87$\\
Gramian direction & $+2.69\times10^{-3}$ & $0.88$ & $0.80$ & $0.69$\\
Gramian direction, reversed & $-2.82\times10^{-3}$ & $1.15$ & $1.27$ & $1.47$\\
random orientation, Gramian & --- & $1.012$ & $1.032$ & $1.081$\\\midrule
for scale: flux--shear split in McCormick's model (Sec.~\ref{sec:mcc}) & $+3.1\times10^{-4}$ & $0.985$ & $0.974$ & $0.958$\\\bottomrule\end{tabular}\end{center}
\grade{estimate: standard Gaussian abundance applied to the measured $\delta_c$ shifts; script \texttt{abundance.py}}.
\begin{itemize}
\item A hidden angular record that is invisible to every fluid and metric quantity at decoupling changes the number of black holes by 5--47\% at the planted amplitude, with a sign set by its orientation. That is about three to eleven times what the flux--shear split does in McCormick's physical model.
\item Because $\beta$ is convex in $\delta_c$, a record of random orientation does not average away: it still \emph{raises} the abundance, by 0.3--1.4\% for $\ell=4$, 0.6--3.8\% for $\ell=3$ and 1.2--8.1\% along the Gramian's leading direction, rarer black holes gaining more.
\end{itemize}
Two cautions apply. McCormick sets abundance with the linear compaction $C_l$, not the compaction peak $\delta_c$ used here, so these are first estimates. And the amplitudes are the engineered half-margin amplitudes of Section~\ref{sec:door}; only that amplitude was run, so how the shift, and with it every number in the table, scales with the size of the record is untested.

\paragraph{Where in mass.} With the horizon mass $M_H(T)\simeq1.5\times10^5\,M_\odot\,(T/\mathrm{MeV})^{-2}$ for $g_*=10.75$, and $\epsilon\propto a\propto1/T$ in radiation, a perturbation whose collisions switch off at $\epsilon_{\rm dec}$ has mass $M\simeq\epsilon_{\rm dec}^2M_H(T_{\rm dec})$. For $T_{\rm dec}\approx1$\,MeV the values $\epsilon_{\rm dec}=0.5$, $1$, $2$, $3$, $4$, $8$ correspond to about $4\times10^4$, $1.5\times10^5$, $6\times10^5$, $1.3\times10^6$, $2.4\times10^6$ and $10^7\,M_\odot$ \grade{estimate}. The door runs of Section~\ref{sec:door} ($\epsilon_{\rm dec}=2$) sit at about $6\times10^5\,M_\odot$, where McCormick's step in the mass function begins. Supermassive black-hole seeds sit in this range, so the effect lands where the observational question is live~\cite{Dave}.

\paragraph{A universe without a first moment.} The standard reading starts the radiation era with the hidden moments empty. McCormick's runs and the memory curves of Section~\ref{sec:memory} begin that way. Then the door has nothing to carry, and the threshold is fixed by the history of the rate alone.

The author's reading is different~\cite{Atlas}. The universe is a self-recurring transport system: closed but growing, relaxing toward full relaxation and then letting go again. Its beginning is ``not a Big Bang of energy, but a first spectral event that forces the atlas to remember''; its first and lightest carrier of the record is light. In that reading the neutrinos do not reach decoupling blank. Full relaxation wipes everything but the kept sector, and the kept sector is what carries the record across. When the gas lets go at the wall, the record is released into the streaming moments. The door of Section~\ref{sec:door} then writes it into the black-hole threshold.

The two readings make different predictions. Both are stated so that observation can decide them:
\begin{enumerate}
\item \textbf{Blank start (standard):} no hidden record at decoupling. The black-hole mass function across $10^5$--$10^7\,M_\odot$ follows McCormick's step with no extra shift or spread.
\item \textbf{Recurrent start (author):} a hidden record at decoupling. If it is coherent in orientation over the scales that form black holes, the abundance shifts by tens of per cent across the neutrino step, up or down according to its parity. If it is incoherent, the step is raised by a few per cent and its edge broadened.
\end{enumerate}
\grade{H: a hypothesis with a stated test}. Cyclic cosmologies exist in other forms. Penrose's conformal cyclic cosmology carries information across aeons through conformal geometry~\cite{CCC}, and Steinhardt and Turok's model through colliding branes~\cite{ST}. The reading here differs in its carrier: a kept transport sector and a hidden angular record, with a measured channel (the door) by which the record reaches matter.

\paragraph{The neutrinos sit on the wall.} The same wall governs the neutrinos' masses. In the author's link coordinate $\Gamma=(z-1)/(z+1)$, two equal perpendicular steps turn parabolic at $\Gamma^2=\tfrac12$, i.e.\ $z=3+2\sqrt2=5.8284$~\cite{SMball}. The measured neutrino splittings (NuFIT~6.0) give, for normal ordering with $m_1=0$, $m_3/m_2=5.8165\pm0.079$, which is $-0.15\sigma$ from the wall, and $\Delta m^2_{31}/\Delta m^2_{21}=33.83$ against $(1+\sqrt2)^4=33.97$. With the first JUNO measurement $\Delta m^2_{21}=(7.50\pm0.12)\times10^{-5}\,$eV$^2$ and NuFIT~6.1's $\Delta m^2_{31}=(2.529\pm0.021)\times10^{-3}\,$eV$^2$, the inputs Alikhanov uses~\cite{Alik}, the ratio is $33.72$ and $m_3/m_2=5.807\pm0.052$, $-0.41\sigma$ from the wall. The relation itself was proposed by Alikhanov~\cite{Alik}, who states it for normal ordering with $m_1=0$. What is the author's is its reading as the wall $\Gamma^2=\tfrac12$. The charged fermions' generation steps all lie past the wall, and the neutrinos' heaviest step lies on it. The particles whose decoupling this paper studies are the particles at the wall.
\begin{itemize}
\item \textbf{Prediction:} normal ordering, and a lightest neutrino mass below about $1.5$\,meV, where the miss on the current data reaches $-2\sigma$ (about $2$\,meV on NuFIT~6.0 alone); inverted ordering misses by far.
\item \textbf{Who can test it:} JUNO, the long-baseline experiments and cosmological mass bounds.
\end{itemize}
\grade{checked numbers; reading; prediction testable by others}.

\paragraph{What the rest of the paper supplies.}
\begin{itemize}
\item Section~\ref{sec:cut}: the wall and the cut.
\item Sections~\ref{sec:runs}--\ref{sec:memory}: that the history, not only its end point, sets the share of the neutrino effect.
\item Section~\ref{sec:witness}: where the record lives.
\item Section~\ref{sec:link}: that the threshold reads the direction of the flux--shear mismatch.
\item Section~\ref{sec:gramian}: which hidden directions matter most.
\item Section~\ref{sec:door}: the measured threshold shifts used above.
\item Section~\ref{sec:mcc}: what McCormick's physical collision model does.
\end{itemize}
The abundance estimates above follow from those measurements in one line each.


\section{The collision operator as a cut}\label{sec:cut}

For one Fourier mode $k$, the linearised neutrino intensity in McCormick's relaxation model evolves under
\begin{equation}
\partial_t f=L f,\qquad L=A+D,\qquad A=-ik\mu,\qquad D=-\kappa(I-P),
\end{equation}
where $\mu$ is the cosine of the angle to the wave vector and $P$ projects onto the isotropic part. Free streaming $A$ is anti-Hermitian. Collisions $D$ are Hermitian, conserve the isotropic part (the neutrino energy in the plasma frame) and remove the flux, which goes to the plasma. In the normalised Legendre basis $|\ell\rangle$, streaming couples only neighbours, $\langle\ell+1|\mu|\ell\rangle=c_\ell=(\ell+1)/\sqrt{(2\ell+1)(2\ell+3)}$, so the operator lives on a chain $\ell=0,1,2,\dots$ (Figure~\ref{fig:chain}).

\paragraph{Non-normality is connection energy.} For Hermitian $D$ and anti-Hermitian $A$, $[L,L^\ast]=2[A,D]$, and for diagonal $D$, $[A,D]_{ij}=A_{ij}(D_j-D_i)$. Hence
\begin{equation}
\tfrac18\lVert[L,L^\ast]\rVert_F^2=\sum_{\text{links }\ell\to\ell+1}k^2c_\ell^2\,(D_{\ell+1}-D_\ell)^2 .
\end{equation}
This is the scalar case of the identity in~\cite{CE}: the self-commutator energy of a transport operator is the energy of its Hermitian field failing to be parallel along the links \grade{derived; Appendix A.1}. With one rate, only the link $0\to1$ has a jump, so
\begin{equation}
\tfrac18\lVert[L,L^\ast]\rVert_F^2=\frac{\kappa^2k^2}{3}\qquad\text{\grade{derived; checked to $10^{-12}$, I1}}.
\end{equation}
All the non-normality of the single-rate model lies on the link between density and flux, the same channel as the momentum drag $\kappa q$ in McCormick's plasma equation~\cite[Eq.~(6)]{McC}.

\paragraph{Beyond relaxation time.} Let the flux relax at $\kappa_1$ (momentum transfer) and shear and higher moments at $\kappa_2$. Then two links carry jumps:
\begin{equation}
\tfrac18\lVert[L,L^\ast]\rVert_F^2=\frac{k^2\kappa_1^2}{3}+\frac{4k^2}{15}(\kappa_2-\kappa_1)^2\qquad\text{\grade{derived; checked, I2}}.
\end{equation}
The energy is stationary in $\kappa_2$ at $\kappa_2=\kappa_1$. If neutrino--neutrino scattering damps shear but not flux, with rates relative to e$^\pm$ scattering of $1.5$ for $\nu_e$ and $7$ for $\nu_{\mu,\tau}$, the shear link carries 64\% and 97.5\% of the non-normality \grade{computed, under those assumed ratios}.

\paragraph{The weight across the cut is the sound speed.} With $H=k\mu$ the weight between kept and streaming parts is $E_B=\lVert PH(I-P)\rVert^2=k^2c_0^2=k^2/3$, so $\sqrt{E_B}=k/\sqrt3=c_s k$, the radiation sound speed~\cite{WB} \grade{derived; I3}. A rate bound built on $E_B$ therefore sees the root-mean-square streaming speed $1/\sqrt3$, not the light-speed front that empties the overdensity~\cite{OSL}.

\paragraph{The kept sector, exactly.} The isotropic part has the closed propagator
\begin{equation}
\langle0|(s-L)^{-1}|0\rangle=\frac{g}{1-\kappa g},\qquad g=\frac1k\arctan\frac{k}{s+\kappa}\qquad\text{\grade{derived; checked, I4}}.
\end{equation}
The two-moment truncation, keeping $\ell=0,1$, has eigenvalues $-\kappa/2\pm\sqrt{\kappa^2/4-k^2/3}$, which merge at $\kappa=2k/\sqrt3$ \grade{derived; I5}. This exceptional point belongs to the truncation. The full chain has its own wall. The propagator is the charge-diffusion channel of the relaxation-time (BGK) equation, and Romatschke found its hydrodynamic pole and the wall~\cite[Eq.~(37)]{Rom}; in the present variables it reads as follows. A mode is a pole, $\kappa g=1$, so $\arctan[k/(s+\kappa)]=k/\kappa$ and
\begin{equation}
s=k\cot\frac{k}{\kappa}-\kappa,
\end{equation}
and this pole exists only for $\kappa>\kappa_c=2k/\pi\approx0.6366\,k$ \grade{derived; Appendix A.3a}. The mode is real all the way down; at $\kappa_c$ it reaches the branch cut of $g$, the segment $\mathrm{Re}\,s=-\kappa$, $|\mathrm{Im}\,s|\le k$, at its centre $s=-\kappa$, and merges into the continuous spectrum: Romatschke's hydrodynamic onset transition, $\tau_D=\pi/(2k)$~\cite{Rom,Hu}. Read here, it is the edge of the neutrino kept sector. With $\phi=k/\kappa$ the mode's non-orthogonality factor is
\begin{equation}
K=\frac{4\phi^2}{\sin^2 2\phi},
\end{equation}
which is $1$ for a fully coupled gas and diverges at the onset $\phi=\pi/2$ \grade{derived; Appendix A.3a}: the normalised overlap of the kept mode with its own left eigenvector, $K^{-1/2}=\sin2\phi/2\phi$, falls to zero, and its spectral projector becomes infinitely ill-conditioned as the mode meets the cut. A truncated chain approaches the onset from above: its least-damped mode turns real at $\kappa/k=0.665$, $0.645$ and $0.639$ with $50$, $200$ and $800$ moments, and at $400$ moments its pole and $K$ match the closed forms to ten and six digits (Figure~\ref{fig:onset}) \grade{computed; Appendix B.26}.



\begin{figure}[t]\centering\includegraphics[width=0.95\textwidth]{figs/fig2_chain.png}
\caption{The neutrino moments form a chain. Streaming links each moment to its neighbours; collisions put a rate on each. Non-normality sits only where neighbouring rates differ: on the density--flux link ($k^2\kappa_1^2/3$) and, if shear relaxes at its own rate, on the flux--shear link. Every other link carries none.}\label{fig:chain}\end{figure}

\begin{figure}[t]\centering\includegraphics[width=0.95\textwidth]{figs/fig_v7_onset.png}
\caption{The wall of the full chain. Left: the kept-sector pole $s=k\cot(k/\kappa)-\kappa$ (line) and the least-damped eigenvalue of a 400-moment chain (dots); the pole meets the branch cut at its centre $s=-\kappa$ (dashed) at $\kappa=2k/\pi$. Right: its non-orthogonality $K=4\phi^2/\sin^22\phi$, with $\phi=k/\kappa$, diverging at the onset. Inset numbers: the truncated chain's onset at 50, 200 and 800 moments, closing on $2/\pi$ from above.}\label{fig:onset}\end{figure}

\section{New runs of McCormick's code}\label{sec:runs}

The code~\cite{McCcode} was run unchanged except for one copied file in which the decoupling exponent $n$ in $\kappa=H(\epsilon_{\rm dec}/\epsilon)^n$ is an input. Pure radiation, $R_\nu=0.409$, $N=400$ cells, $N_\mu=32$ angle bins, $\epsilon_0=0.1$, as in his Table~II. A reproduction run at $\delta=0.5109$ with $n=3$, $\epsilon_{\rm dec}=2$ forms a black hole, inside his published bracket $(0.51074,0.51109)$ \grade{measured; Appendix B.4}. For the steep history, $\kappa/H$ was capped at $5\times10^5$, the stiffness his own $\epsilon_{\rm dec}=8$ runs reach at $\epsilon_0$. At that cap deviations from isotropy are of relative order $H/\kappa\le2\times10^{-6}$, the same estimate McCormick uses for his switch to transport~\cite[Sec.~III]{McC}.

For every test below each prediction and its kill condition were written before the runs (Appendix~B.3). The threshold search runs two points per round; a run that neither collapses nor disperses by the end time, or that stops on an error, never moves the bracket, is retried once with a longer end time and is reported if still undecided. No point in any run below was undecided. These runs use McCormick's own end-state rule, in which a numerical breakdown after $2M/R>0.95$ counts as collapse; Section~\ref{sec:link} tightens that rule to a detected horizon and audits every bracket edge under it.

\paragraph{Test A: a profile that collapses late.} The broad profile $\alpha=0.3$ is decided near $\epsilon\approx7$--$9$, against $\epsilon\approx4$ for the Gaussian. Its thresholds are $\delta_c^{\rm fluid}=0.44805$ and $\delta_c^{\rm inst}=0.47747$~\cite{McCcode}. With $n=3$:
\begin{center}\begin{tabular}{lccc}\toprule
& $\delta_c$ & $G$ (broad) & $G$ (Gaussian, \cite{McC}) \\\midrule
$\epsilon_{\rm dec}=4$ & $0.45337\pm0.00013$ & $0.181\pm0.004$ & $0.190$\\
$\epsilon_{\rm dec}=8$ & $0.44938\pm0.00012$ & $0.045\pm0.004$ & $0.063$\\\bottomrule\end{tabular}\end{center}
The broad profile keeps the same share of the neutrino effect as the Gaussian or less (29\% less at $\epsilon_{\rm dec}=8$), although it collapses about twice as late (Figure~\ref{fig:clock}) \grade{measured}. A clock tied to the collapse time would have raised $G$; it fell. Because the test changes the profile shape together with the collapse time, it shows that the collapse time alone does not set the clock; it does not by itself separate the two, and it does not fix which earlier epoch does.

\paragraph{Tests B, C and D: a steeper decoupling history.} Gaussian profile, $\kappa/H=(\epsilon_{\rm dec}/\epsilon)^6$:
\begin{center}\begin{tabular}{lccc}\toprule
& $\delta_c$ & $G$ ($n=6$) & $G$ ($n=3$, \cite{McC}) \\\midrule
C: $\epsilon_{\rm dec}=1$ & $0.51756\pm0.00019$ & $0.648\pm0.006$ & $0.683$\\
B: $\epsilon_{\rm dec}=2$ & $0.50743\pm0.00011$ & $0.317\pm0.004$ & $0.431$\\
D: $\epsilon_{\rm dec}=3$ & $0.50245\pm0.00019$ & $0.154\pm0.006$ & $\approx0.30$\\\bottomrule\end{tabular}\end{center}
\grade{measured}. The $n=3$ value at $\epsilon_{\rm dec}=3$ lies between McCormick's $0.431$ and $0.190$ at 2 and 4. A neutrino population that decouples more abruptly, but by the same nominal time, keeps less of its effect, and the gap widens with $\epsilon_{\rm dec}$. The history matters, not only its end point.

\paragraph{Numerical controls.} Test B was repeated with the rate cap raised tenfold, to $5\times10^6$, and on a finer grid, $N=600$ cells and $N_\mu=48$ angle bins. The thresholds are $\delta_c=0.50740\pm0.00011$ and $0.50751\pm0.00006$, shifts of $-0.00003$ and $+0.00008$ against the pre-registered tolerance of $0.0005$ \grade{measured}. The cap acts only while the neutrinos are tightly coupled: the uncapped rate starts at $(\epsilon_{\rm dec}/\epsilon_0)^6$ and falls below the cap at $\epsilon\approx0.11$, $0.22$ and $0.34$ for $\epsilon_{\rm dec}=1$, 2 and 3. The cap control was run at $\epsilon_{\rm dec}=2$ only; at $\epsilon_{\rm dec}=1$ and 3 it is untested.

\begin{figure}[t]\centering\includegraphics[width=0.9\textwidth]{figs/fig3_clock.png}
\caption{Test A. The broad profile collapses at $\epsilon\approx7$--$9$, yet keeps the same share $G$ of the free-streaming shift as the Gaussian. A clock tied to the collapse time would have given the dotted values.}\label{fig:clock}\end{figure}

\section{What does not set the threshold}\label{sec:not}

Each reading below was tested against numbers. Only the last two were tested against runs set up for them in advance; the first three were tested against McCormick's published curve.

\begin{enumerate}
\item \textbf{The first maximum of linear theory.} The linear density contrast at its first maximum ($\epsilon\approx1.7$) reproduces the instantaneous shift (centre factor $0.944$, as McCormick finds). Under gradual decoupling the share it keeps is 0.86, 0.59, 0.19, $-0.05$, $-0.02$ and 0.00, a largest error of $0.50$. Decoupling after $\epsilon\approx1.7$ is nearly invisible to it, yet still moves the nonlinear threshold \grade{computed}.
\item \textbf{A sharp wall, and rate ratios.} Treat the two-moment exceptional point, crossed at $\epsilon_{\rm EP}=0.811\,\epsilon_{\rm dec}^{3/4}$ for $kr_m=2$, as a wall, or use the parameter-free escape fractions $k/(k+\kappa)$, $e^{-\kappa/k}$ and $1-e^{-k/\kappa}$. Averaging each over two windows, with two weightings and two scales, gives 32 combinations. The wall variants miss by $0.14$ or more. With the window starting at horizon entry, every combination misses by $0.17$ or more. One combination, $e^{-\kappa/k}$ averaged from $\epsilon_0$, comes within $0.077$; selected from 32, it is not counted. On Test B, which it never saw, it predicts $0.513$ against $0.317$ \grade{computed; measured}.
\item \textbf{Diffusive escape.} Diffusion with $D=1/(3\sigma)$ from the conformal rate $\sigma=\epsilon_{\rm dec}^3/\tau^4$, combined in series with streaming at the collapse time $\epsilon^\ast$, gives $G=1/[1+\sqrt{15}(\epsilon_{\rm dec}/\epsilon^\ast)^{3/2}]$. With $\epsilon^\ast=4$ it matches McCormick's curve to $0.02$, but it was built after seeing that curve, and a free fit of its exponent gives $1.61\pm0.03$ rather than $3/2$. Test A with $\epsilon^\ast=7$--$9$ predicted $0.37$--$0.47$ and $0.17$--$0.24$, and Test B predicted $0.62$. Both are excluded \grade{measured}.
\item \textbf{History-blind.} If only $\epsilon_{\rm dec}$ mattered, Tests B and D would give $0.431$ and about $0.30$. They give $0.317$ and $0.154$; Test D misses by $0.146$ against a kill line of $0.08$ \grade{measured}.
\end{enumerate}

\section{What survives: the memory of the kept sector}\label{sec:memory}

The reading that survives treats the isotropic part as a kept sector and asks how much of it the decoupling operator keeps~\cite{WB}. Start a mode in $|0\rangle$ at $\tau_0=0.02$ and evolve it under $\partial_\tau v=-(ik\mu+a\kappa(\tau)(I-P))v$, with $a\kappa=\epsilon_{\rm dec}^n/\tau^{n+1}$. Let $A(\tau)$ be the population left in $|0\rangle$, averaged over the Gaussian profile with weights $k^2e^{-k^2/2}$. Define
\begin{equation}
g(\epsilon^\ast)=\frac{1-A_{\epsilon_{\rm dec}}(\epsilon^\ast)}{1-A_{\rm inst}(\epsilon^\ast)},
\end{equation}
the share of the kept sector's loss that has happened by the readout time $\epsilon^\ast$.

With no fitted number, read at the first compaction maximum $\epsilon^\ast=1.7$, it gives 0.941, 0.860, 0.671, 0.344, 0.087 and 0.013 against McCormick's $G$. That is within $0.02$ for $\epsilon_{\rm dec}\le1$, with a largest error of $0.103$, just outside the $0.10$ bar set for it in advance \grade{computed}. McCormick's own curve passes $G=\tfrac12$ at $\epsilon_{\rm dec}\approx1.67$, within 3\% of the first maximum: the neutrinos must be decoupled by then for half the effect to survive.

The late shortfall says the collapse is still listening after the first maximum. Mixing the two epochs McCormick's numbers single out, $g_2=w\,g(1.7)+(1-w)\,g(4.0)$, with one weight fitted to his curve, $w=0.89$, gives a largest error of $0.022$ \grade{fit}. The form is not unique: with the scale fixed at $kr_m=2$ the single-epoch reading misses by $0.093$, and the best single readout epoch, $\epsilon^\ast=1.87$, by $0.065$. Its weight does not carry over between weightings: the same $w$ with the scale fixed at $kr_m=2$ misses by $0.115$, outside the $0.10$ bar declared for that check \grade{computed}. The form is a fitted shape with a stated sensitivity, not a derived law.

The test is what that fitted form says about histories it never saw. All its $n=6$ values were computed in one file dated 13:47 on 30 October, before any $n=6$ run, and were not refitted afterwards:
\begin{center}\begin{tabular}{lccccc}\toprule
& predicted & measured & miss & history-blind & miss\\\midrule
C: $\epsilon_{\rm dec}=1$ & $0.663$ & $0.648$ & $0.014$ & $0.683$ & $0.035$\\
B: $\epsilon_{\rm dec}=2$ & $0.246$ & $0.317$ & $0.071$ & $0.431$ & $0.114$\\
D: $\epsilon_{\rm dec}=3$ & $0.118$ & $0.154$ & $0.037$ & $\approx0.30$ & $0.146$\\\bottomrule\end{tabular}\end{center}
Every miss of the memory form lies inside the kill line of $0.08$ set before its run; the history-blind rule fails at $\epsilon_{\rm dec}=3$ (Figures~\ref{fig:G} and~\ref{fig:kill}) \grade{measured}. At $\epsilon_{\rm dec}=2$ and 3 the form falls below the measurement: it counts slightly too much of the early erasure. Test C was set where the two rules nearly agree, as a survival check; Test D was set where they separate. A readout window set by the neutrino flux and the enclosed mass between first maximum and apparent horizon, as McCormick suggests~\cite[Sec.~IV\,C]{McC}, is the natural next object. It is untested.

\begin{figure}[t]\centering\includegraphics[width=0.95\textwidth]{figs/fig4_tests.png}
\caption{The surviving share $G$ of the free-streaming shift against decoupling time. Black dots: McCormick's Gaussian runs with $\kappa/H=(\epsilon_{\rm dec}/\epsilon)^3$; the grey dashed line through them is what a history-blind rule gives for any history. Blue: the kept-sector memory with one weight fitted to the black dots. Orange line: the same form for $\kappa/H=(\epsilon_{\rm dec}/\epsilon)^6$, computed before any such run and not refitted. Orange diamonds: the new $n=6$ runs C, B and D. Open squares: the broad profile (Test A).}\label{fig:G}\end{figure}

\begin{figure}[t]\centering\includegraphics[width=\textwidth]{figs/fig5_killlines.png}
\caption{The three steep-history tests against their kill lines. The shaded band is $\pm0.08$ around each measurement. Orange: the memory form's prediction, made before the run. Grey circle: the history-blind rule. At $\epsilon_{\rm dec}=3$ the history-blind rule falls outside the band.}\label{fig:kill}\end{figure}

\section{Where the memory sits}\label{sec:witness}

A fluid closure carries the density, the flux and the anisotropic stress, the moments $\ell=0,1,2$. The memory form of Section~\ref{sec:memory} says the threshold depends on how the kept sector was emptied. The question here is whether that history can sit in moments a closure does not carry.

\paragraph{Exact elimination of the hidden moments.} Split $v=(a,h)$ with $a=v_0$ and $h=(v_1,v_2,\dots)$. In the relaxation model every hidden moment is damped at the same rate, so the collision term on the hidden block is $\sigma(\tau)$ times the identity and commutes with streaming. The hidden propagator is therefore exact in closed form,
\begin{equation}
U_Q(t,r)=e^{-[\Sigma(t)-\Sigma(r)]}\,e^{-ik(t-r)M_{QQ}},\qquad \Sigma(\tau)=-\frac{\epsilon_{\rm dec}^{\,n}}{n\tau^n},\quad \Sigma'=\sigma,
\end{equation}
and eliminating $h$ gives one equation for the kept amplitude,
\begin{equation}
\dot a(t)=S(t)-k^2\!\int_s^t K(t,r)\,a(r)\,dr,\qquad K=c_0^2\,[U_Q(t,r)]_{11},\qquad S=-ikc_0\,[U_Q(t,s)h(s)]_1
\end{equation}
\grade{derived; Appendix A.5}. The kernel $K$ is the memory; the source $S$ is everything the hidden moments held at time $s$.

\paragraph{The matched pair.} Take two states with the same $(v_0,v_1,v_2)=(1,0.2,0.1)$, differing only by $\pm0.05$ in $v_3$ (or $v_4$), scaled by one constant to unit norm. Both are positive angular distributions ($f(\mu)\ge0.50$). Their kept amplitudes start equal, so their difference obeys the equation above driven only by $S$ of the hidden difference. Both are evolved under the same history, $\epsilon_{\rm dec}=0.5$, with the planting time $\tau_s$ scanned across decoupling. The readout is the profile-weighted change in $(v_0,v_1,v_2)$, relative to its size, $R_{\rm ret}$; the witness is on if $R_{\rm ret}>10^{-3}$. Every on/off call was written down before the run, from the survival factor of the hidden block, $\exp[-\epsilon_{\rm dec}^n/(n\tau_s^n)]$, and the leakage $(k/\sigma)^\ell$ through damped rungs (Appendix~B.14).

\begin{center}\begin{tabular}{lcccc}\toprule
$\tau_s/\epsilon_{\rm dec}$ & $n=3$, $\ell=3$ & $n=3$, $\ell=4$ & $n=6$, $\ell=3$ & $n=6$, $\ell=4$\\\midrule
0.1 & $7\times10^{-15}$ & $5\times10^{-19}$ & $8\times10^{-24}$ & $5\times10^{-31}$\\
0.25 & $6\times10^{-10}$ & $6\times10^{-11}$ & $2\times10^{-15}$ & $7\times10^{-20}$\\
0.5 & $\mathbf{8.6\times10^{-3}}$ & $\mathbf{7.9\times10^{-3}}$ & $3\times10^{-6}$ & $3\times10^{-6}$\\
1 & $\mathbf{0.081}$ & $\mathbf{0.074}$ & $\mathbf{0.094}$ & $\mathbf{0.086}$\\
2 & $\mathbf{0.095}$ & $\mathbf{0.085}$ & $\mathbf{0.098}$ & $\mathbf{0.088}$\\\bottomrule\end{tabular}\end{center}
$R_{\rm ret}$ at $\tau=4$; bold entries are on \grade{computed}. Every call came out as predicted. The survival law is quantitative: for $n=3$ the ratio of the entries at $0.5$ and $1$ is $0.106$ against $e^{-8/3}/e^{-1/3}=0.097$, and for $n=6$ at $0.5$ the law gives $1.4\times10^{-6}$ against $1.6\times10^{-6}$ at $\tau=1.7$. Planted in the retained moments instead, deep in coupling the order is $\ell=1>2>3>4$, each rung costing about $10^{-4}\approx k/\sigma$; after decoupling all four are of the same size. A hidden difference set at decoupling changes the later density, flux and stress as much as a visible one (Figure~\ref{fig:door}).

\paragraph{Checks.} Doubling the number of moments, $60\to120$, and tightening the solver tolerances change every entry by at most $1.4\times10^{-5}$ of its value. Solving the elimination equation from the closed-form $U_Q$ alone, with no solve of the full system, reproduces the full difference to $3\times10^{-7}$--$7\times10^{-7}$, falling fourfold with each halving of the step; leaving out the kernel misses by 19--234\% \grade{computed; Appendix B.14}.

\paragraph{A speed limit on the kept population.} Collisions conserve $|0\rangle$ and remove the rest, so $\lVert v\rVert$ never grows, and streaming links $|0\rangle$ only to $|1\rangle$, with weight $kc_0=\sqrt{E_B}$. Hence for $A=|v_0|^2$
\begin{equation}
\dot A^2\le4E_B\,A(1-A),\qquad E_B=k^2/3
\end{equation}
\grade{derived; Appendix A.6; Lean}. On every trajectory computed here the largest value of the ratio is $0.87$ \grade{computed}. The bound limits how fast the kept sector can be emptied; it does not set the threshold.

\paragraph{Reading.} For structure set while the neutrinos are tightly coupled, a closure on density, flux and stress loses almost nothing: planted at a quarter of the decoupling time or earlier, the hidden difference returns at $10^{-9}$ or less in the table above. Structure set at or after decoupling survives in the finer angular moments and streams back into the visible ones by several per cent. The memory McCormick's runs record therefore has a place to live that no fluid closure contains. How far such structure moves the nonlinear threshold is measured for matched, planted perturbations in Section~\ref{sec:door}. Whether realistic initial data carry it is untested: the memory curves of Section~\ref{sec:memory} start with the hidden moments empty.

\begin{figure}[t]\centering\includegraphics[width=\textwidth]{figs/fig6_door.png}
\caption{The door. Strength of the matched-pair witness, $\log_{10}R_{\rm ret}$ at $\tau=4$, against the time the hidden difference was planted (in units of $\epsilon_{\rm dec}$) and the moment that held it. Planted deep in the coupled era the difference is erased, down to $10^{-30}$; planted at or after decoupling it returns at the $10^{-1}$ level. The red plane is the witness threshold $10^{-3}$. The steeper history ($n=6$) closes the door more sharply.}\label{fig:door}\end{figure}

\section{The shear sets the threshold}\label{sec:link}

Section~\ref{sec:cut} put all the non-normality of a single-rate collision term on one link. A collision term with two rates has two links that carry it: density to flux, with jump $\kappa_1$, and flux to shear, with jump $\kappa_2-\kappa_1$. The connection energy of the second link, $\tfrac{4k^2}{15}(\kappa_2-\kappa_1)^2$, is stationary at $\kappa_2=\kappa_1$, while the exact kept response is not (Appendix~A.7). Which does the nonlinear threshold follow, and which moment carries what it follows? This section answers both with McCormick's code, changed in one place, and then asks in which coordinate the answer is a law.

\paragraph{The change.} In a copy of \texttt{ms\_nuc.py} the collision term $\kappa(\varepsilon/4\pi-I)$ becomes $\kappa_2(\varepsilon/4\pi-I)+(\kappa_2-\kappa_1)d_\mu$, where $d_\mu$ is the dipole part of the intensity, projected on the same angle bins that define the flux (Appendix~A.12). The dipole then relaxes at $\kappa_1$ and every moment $\ell\ge2$ at $\kappa_2$; the neutrino energy is untouched, and the momentum handed to the plasma is still $\kappa_1q$. The stiff time-step cap uses the larger rate. Write $r=\kappa_2/\kappa_1$, and keep $\kappa_1=\kappa=H(\epsilon_{\rm dec}/\epsilon)^3$ with $\epsilon_{\rm dec}=2$, the Gaussian profile, $R_\nu=0.409$, $N=400$, $N_\mu=32$. At $r=1$ the changed code reproduces the unchanged one: $\delta=0.5109$ forms a black hole at $\epsilon=6.61$ in both \grade{measured}.

\paragraph{Admissible ratios.} The split collision operator is a Markov generator for $\tfrac34\le r\le\tfrac32$ (all moments $\ell\ge2$ at $\kappa_2$) and $0.6\le r\le1.2$ (shear alone at $\kappa_2$), the conservative continuum windows; on the $32$ angle bins the exact windows are slightly wider, $0.725\le r\le1.55$ and $0.56\le r\le1.24$, and outside them the operator has a negative off-diagonal entry and can, in principle, drive a nonnegative intensity negative. The pairs $r=1.2$, $1/1.2$, $1.25$ and $0.8$ are inside the first window. The pairs $r=1/2$ and $2$ and the extreme pairs ($r=1/16,1/4,4,16$) are formal interventions on the solver: they are kept because they fix the sign and the size of the effect, and the minimum intensity was logged through every run of this section; it stayed positive in all of them except one run at $r=16$, where it reached $-2.4\times10^{-10}$ \grade{measured}.

\paragraph{How the searches are run.} Each threshold is found by trisection: two trial amplitudes per round, the bracket shrinking to its middle third. A run moves the bracket only if it ends in a detected horizon ($2M/R\ge1$ with $U<0$) or disperses. A run that stops on a numerical error, or that breaks down after $2M/R>0.95$ without a detected horizon, is repeated once with the Courant numbers lowered from $0.4$ to $0.3$; if it fails again it is reported and never moves the bracket. Every decisive bracket edge of every threshold below was rerun under this rule: each upper edge is a detected horizon and each lower edge a dispersal \grade{measured; Appendix~B.15}.

\paragraph{The link as a reflecting step.} In the author's coordinate a link between rates $d_i$ and $d_j$ is one reflecting step, with Smith coordinate $\Gamma_{ij}=(d_i-d_j)/(d_i+d_j)=\tanh(\tfrac12\ln d_i/d_j)$ and rapidity $\tfrac12\ln(d_i/d_j)$; exchanging the ends negates $\Gamma$ (Appendix~A.8), and a matched link carries no connection energy~\cite{CE}. For the flux--shear link, $\Gamma_{12}=(r-1)/(r+1)$. The surviving share $G$ is itself a point on the disk, $\Gamma_G=1-2G$, with offset $\tau_G=\ln[(1-G)/G]$. The question becomes: how does $\tau_G$ move with the link's rapidity?

\paragraph{Three readings, written down before the runs.} (a) The frozen memory form of Section~\ref{sec:memory}, with the collision diagonal $(0,1,r,r,\dots)$, predicts almost no effect: $\Delta G=-0.004$ at $r=1/4$ and $+0.003$ at $r=4$. (b) If the connection energy controlled the threshold, the two shifts would share a sign, since the energy is even about $r=1$. (c) If the threshold read the link's reflection coefficient, the shifts would be odd in $\Gamma_{12}$: equal and opposite.

\paragraph{Eleven thresholds.} The reference thresholds are McCormick's, $\delta_c^{\rm fluid}=0.49773$ and $\delta_c^{\rm inst}=0.52832$, and $G=(\delta_c-\delta_c^{\rm fluid})/(\delta_c^{\rm inst}-\delta_c^{\rm fluid})$; the point $r=1$ is measured here with the changed code. These eleven runs hold the collision rate fixed through each time step, as in version~5; the stage-correct reruns of the admissible points are given below:
\begin{center}\begin{tabular}{rrcccc}\toprule
$r$ & $\Gamma_{12}$ & $\tfrac12\ln r$ & $\delta_c$ & $G$ & $\tau_G-\tau_G(1)$\\\midrule
$1/16$ & $-0.882$ & $-1.386$ & $0.521485\pm0.000038$ & $0.777$ & $-1.547$\\
$1/4$ & $-0.600$ & $-0.693$ & $0.516498\pm0.000038$ & $0.614$ & $-0.763$\\
$1/2$ & $-0.333$ & $-0.347$ & $0.513656\pm0.000010$ & $0.521$ & $-0.384$\\
$1/1.2$ & $-0.091$ & $-0.091$ & $0.5115145\pm0.0000003$ & $0.451$ & $-0.103$\\
$0.8$ & $-0.111$ & $-0.112$ & $0.5116957\pm0.0000004$ & $0.457$ & $-0.127$\\
$1$ & $0$ & $0$ & $0.510741\pm0.000019$ & $0.425$ & $0$\\
$1.2$ & $+0.091$ & $+0.091$ & $0.5100149\pm0.0000005$ & $0.402$ & $+0.098$\\
$1.25$ & $+0.111$ & $+0.112$ & $0.5098407\pm0.0000002$ & $0.396$ & $+0.122$\\
$2$ & $+0.333$ & $+0.347$ & $0.507973\pm0.000010$ & $0.335$ & $+0.385$\\
$4$ & $+0.600$ & $+0.693$ & $0.505442\pm0.000038$ & $0.252$ & $+0.786$\\
$16$ & $+0.882$ & $+1.386$ & $0.501673\pm0.000025$ & $0.129$ & $+1.610$\\\bottomrule\end{tabular}\end{center}
\grade{measured}. At $r=1/4$ and $4$ the shifts from the measured $G(1)=0.425$ are $+0.188$ and $-0.173$, so reading (a) fails by a factor of forty and with the wrong sign at $r=4$, and reading (b) fails on the sign. Reading (c) holds in the offset, where the law is stated: $-0.763$ and $+0.786$, equal and opposite to $3\%$ (Figure~\ref{fig:link}).

\paragraph{One line through the matched link.} Read in the offset, the eleven points lie close to one line through $r=1$ (residual $0.014$ in $\tau_G$). Each pair $r,1/r$ splits into an odd part, $[\tau_G(r)-\tau_G(1/r)]/\ln r$, the rate at which the offset moves per unit of link rapidity, and an even part. The anchor $\tau_G(1)$ cancels from the odd part; the two reference thresholds do not, because each enters through $\ln(\delta_c-\delta_c^{\rm fluid})$ at two different $\delta_c$ (Appendix~A.9). The ranges below are hard bounds from the four corners of each pair of brackets:
\begin{center}\small\begin{tabular}{lccccc}\toprule
$|\tfrac12\ln r|$ & $0.091$ & $0.112$ & $0.347$ & $0.693$ & $1.386$\\\midrule
odd rate & $1.1006$ & $1.1128$ & $1.1093$ & $1.1177$ & $1.1384$\\
bracket range & $[1.1000,1.1011]$ & $[1.1124,1.1131]$ & $[1.1052,1.1134]$ & $[1.1092,1.1262]$ & $[1.1332,1.1436]$\\
even part & $-0.003$ & $-0.003$ & $+0.001$ & $+0.012$ & $+0.032$\\\bottomrule\end{tabular}\end{center}
\grade{measured}. The rate is not one: a rule in which the offset of $G$ moves by exactly the link rapidity is excluded at every pair. Inside the admissible window it lies at $1.10$--$1.11$; it rises to $1.138$ at the outermost, formal pair, where a small even part appears.

\paragraph{Stage-correct histories.} In the runs above the collision rate $\kappa(\tau)$ was evaluated once per time step and held through the three stages of the Runge--Kutta step. Evaluating it at each stage time is the only change of Appendix~B.19; it moves black-hole formation at $r=1$, $\delta=0.5109$, from $\epsilon=6.6147$ to $6.6176$. The admissible thresholds and the shear-only pair at $r=1.2$ were rerun with it, each by bisection to a bracket width of at most $10^{-6}$ under the same endpoint rule, with predictions and kill lines written first:
\begin{center}\small\begin{tabular}{lccc}\toprule
$r$ & $\delta_c$, stage-correct & $G$ & change\\\midrule
$1$ & $0.51077853$ & $0.4266$ & $+3.8\times10^{-5}$\\
$1.2$ & $0.51000812$ & $0.4014$ & $-6.7\times10^{-6}$\\
$1/1.2$ & $0.51151594$ & $0.4507$ & $+1.4\times10^{-6}$\\
$1.25$ & $0.50984158$ & $0.3959$ & $+8.8\times10^{-7}$\\
$0.8$ & $0.51169233$ & $0.4564$ & $-3.4\times10^{-6}$\\
shear only, $1.2$ & $0.50995463$ & $0.3996$ & $+4.7\times10^{-7}$\\
shear only, $1/1.2$ & $0.51157978$ & $0.4528$ & $+5.0\times10^{-6}$\\\bottomrule\end{tabular}\end{center}
with half-widths $3$--$5\times10^{-7}$ \grade{measured; Appendix~B.19}. The changes are a few $10^{-6}$, but at $|\tfrac12\ln r|\approx0.1$ they move the rates. The inner odd rates become $1.1066$, range $[1.1061,1.1072]$, for $r=1.2$ and $1/1.2$, and $1.1102$, range $[1.1098,1.1106]$, for $r=1.25$ and $0.8$: the two pairs now differ by $0.0036$, against $0.012$ before, so most of the earlier spread between them was the step-held collision rate and not a per-point systematic. The shear-only rate at $r=1.2$ becomes $1.1928$, range $[1.1922,1.1933]$, still $0.086$ above the all-$\ell\ge2$ rate at the same pair. Of the move at $r=1$, $2.2\times10^{-5}$ is the stage fix and $1.6\times10^{-5}$ is the distance from the earlier midpoint at tolerance $10^{-4}$ to the threshold; it moves $G(1)$ to $0.4266$ and the even parts by about $0.003$, and enters no odd rate.

\paragraph{The error the brackets do not show.} The odd rate divides by $\ln r$, so at the inner pairs a shift of $10^{-4}$ in one threshold moves it by $0.06$. Three systematics are larger than the brackets. The fluid reference carries $\pm7.8\times10^{-5}$ in McCormick's archive, which moves every rate by $\pm0.0038$ together. Refining the $r=1.25$ pair in radius, angle and time together ($N=600$, $N_\mu=48$, the Courant numbers and the Hubble step factor halved) raises both thresholds, by $7.5$ and $6.4\times10^{-5}$, to $0.50991687$ and $0.51175643$, and lowers the rate to $1.1020$, range $[1.1016,1.1024]$: a shift of $-0.0082$, which fails the line of $0.005$ set before the run. The angle alone moved the rate by less than $0.001$ at $r=4$ (below), so the shift comes from the radial grid or the time step; which one is untested. With the pair spread of $0.0036$ the budget is $0.0097$ in quadrature, $0.016$ added linearly. The values at the two finer levels, $1.1020$ and $1.1013$ (below), and the whole bracket range of the finest, lie in
\begin{equation}
\text{odd rate}=1.10\pm0.01;
\end{equation}
the first-level values, $1.1066$ and $1.1102$, sit near its upper edge, the second $0.0002$ above it. The $\pm0.01$ is the error budget itself, led by the third-level brackets, and not a margin around a converged number. A third level ($N=900$, $N_\mu=72$, Courant numbers and Hubble step halved again) was run with its predictions and kill line written first. It puts the thresholds at $0.5099481\pm0.0000063$ ($r=1.25$) and $0.5117877\pm0.0000063$ ($r=0.8$), further rises of $3.1\times10^{-5}$ each against $7.5$ and $6.4\times10^{-5}$ for the previous step. The shifts shrank to $0.42$ (range $0.33$--$0.50$) and $0.49$ (range $0.39$--$0.59$) of the previous ones. Cells and angle bins shrink by $1.5$ per level and the Courant-limited time step by $3$. First order in the cell size predicted $0.50996706$ and $0.51179917$, both outside the brackets; second order in the cell size (a factor $0.44$) predicted $0.50995033$ and $0.51178492$, both inside; first order in the time step (a factor $0.33$) predicted $0.50994197$, at the lower edge of the $r=1.25$ bracket, and $0.51177780$, below the $r=0.8$ bracket. So the brackets exclude an error purely first order in the cell size, and at $r=0.8$ one purely first order in the time step; they agree with second order in the cell size and do not exclude a mix of first-order errors in both \grade{measured; Appendix B.27}. The third-level rate is $1.1013$, range $[1.0938,1.1088]$ over the corners of the two brackets: a change of $-0.0007$ from the second level against $-0.0082$ for the previous step, inside the kill line written first, which was a change larger than the previous one. The third-level brackets are $1.25\times10^{-5}$ wide, since runs this close to threshold no longer finish in one compute window, so the change is known only to $\pm0.0075$ and its sign is not resolved; the order read from the three central values, $6.2$, lies outside the range $0.5$ to $4$ set before the runs for an extrapolation, and none is quoted. With the reference shift and the pair spread the budget is $0.0091$ in quadrature, $0.015$ added linearly. Every threshold here starts at $\epsilon_0=0.1$, as McCormick's Table~II; his runs started at $\epsilon_0=0.05$ and $0.025$ move the reference thresholds by about $4\times10^{-4}$~\cite{McCcode}, and how the odd rate depends on the start time is untested. The rate is a measured number with no closed form attached \grade{measured}.

\paragraph{Which moment carries it.} Damp only the shear, $\ell=2$, at $\kappa_2$, and keep the flux and every $\ell\ge3$ at $\kappa_1$: the collision term becomes $\kappa_1(\varepsilon/4\pi-I)+(\kappa_1-\kappa_2)s_\mu$, with $s_\mu$ the $\ell=2$ projection on the same angle bins (Appendix~A.12). Now two links jump, flux to shear by $+\tfrac12\ln r$ and shear to the next rung by $-\tfrac12\ln r$. A rule in which the offset moves by a sum of link-weight ratios times link rapidities then nearly cancels, since $c_1$ and $c_2$ differ by $1.8\%$, and predicts $|\text{rate}|<0.1$; a threshold that reads the shear damping itself predicts a rate near the one above. Both predictions and their kill lines were written before the runs, first at the extreme pair and then, after the admissibility window was found, inside it. At $r=1$ the changed code again reproduces the unchanged one.
\begin{center}\begin{tabular}{lcccc}\toprule
shear only & $\delta_c$ & $G$ & odd rate & all $\ell\ge2$, same pair\\\midrule
$r=1.2$ & $0.5099542\pm0.0000002$ & $0.400$ & \multirow{2}{*}{$1.1895$, range $[1.1892,1.1898]$} & \multirow{2}{*}{$1.1006$}\\
$r=1/1.2$ & $0.5115748\pm0.0000002$ & $0.453$ & & \\\midrule
$r=4$ & $0.5052350\pm0.0000037$ & $0.245$ & \multirow{2}{*}{$1.2038$, range $[1.2030,1.2046]$} & \multirow{2}{*}{$1.1177$}\\
$r=1/4$ & $0.5170942\pm0.0000037$ & $0.633$ & & \\\bottomrule\end{tabular}\end{center}
\grade{measured}. The link-weight reading fails by an order of magnitude at both pairs; the shear reading holds at both. Damping the shear alone moves the threshold more than damping the shear and every higher moment together, by $0.09$ in the rate at either pair: the moments above the shear, relaxed with it, pull the other way. With stage-correct histories the $r=1.2$ pair gives $1.1928$ against $1.1066$ (above).

\paragraph{Angular resolution.} Repeating the $r=4$ and $1/4$ pair with $48$ angle bins instead of $32$, set before the run to pass if the odd rate moved by less than $0.01$ and each threshold by less than $10^{-4}$, gives $\delta_c=0.5054959\pm0.0000082$ and $0.5165576\pm0.0000041$: both thresholds rise, by $5.3$ and $6.0\times10^{-5}$, and the odd rate moves from $1.1177$ to $1.1170$ (range $[1.1156,1.1184]$). The resolution shift is common to $r$ and $1/r$ and cancels from the odd rate to below $0.001$; it moves $G$ itself by about $0.002$ \grade{measured}.

\paragraph{Two small factors, checked.} The reference point for $r=1$ is measured here with the changed code at tolerance $10^{-4}$: $G=0.4253\pm0.0006$, at the lower edge of McCormick's bracket $[0.4253,0.4367]$, whose midpoint $0.431$ would shift every prediction that uses it by about $0.006$. Only the $r>1$ runs shrink the stiff time step, by the factor $r$; rerunning $r=1/4$ with its step shrunk fourfold leaves $\delta_c$ unchanged to $3\times10^{-6}$, so the slight asymmetry between the two sides is not numerical \grade{measured}.

\paragraph{The neutrino fraction.} The kept-sector operator contains no $R_\nu$. With $R_\nu=0.2$, $n=3$, $\epsilon_{\rm dec}=2$ and McCormick's instantaneous threshold at that fraction, $0.51297$, the threshold is $0.50396\pm0.00003$, so $G=0.409\pm0.005$. The memory form had predicted $0.411$; against $0.425$ measured here at $R_\nu=0.409$ the difference is $0.016$ \grade{measured}. The share is nearly, not exactly, independent of how much of the radiation is neutrinos.

\paragraph{Which coordinate the law lives in.} The rapidity $\tfrac12\ln(\kappa_2/\kappa_1)$ is a coordinate on the parameters of the collision operator. In the author's framework~\cite{IQGT} a parameter chart does not generate the geometry of a kept subspace; the intrinsic measure of how far the kept sector moved when $r$ changed is the distance of its state in the moment chain. That distance was computed for the same histories, as the principal angle at the first compaction maximum between the kept-sector trajectory at ratio $r$ and at $r=1$ (Section~\ref{sec:memory}'s chain, profile-averaged), with the readings and kill lines written first: if the threshold read that distance, the all-$\ell\ge2$ and shear-only families would fall on one line in it, and the rise of the odd rate with $|\tfrac12\ln r|$ would flatten. Neither happens: the two families' slopes in the state angle differ by $5\%$ (they differ by $6\%$ in the rapidity), the odd rates spread by $36\%$ against $3.4\%$, and the eleven offsets are far less linear in the angle (residual $0.13$) than in the rapidity ($0.014$). The two collision models move the linear kept-sector state at readout by almost the same angle ($0.0475$ against $0.0467$ at $r=4$) while their thresholds differ by $8\%$. The eigenprojector of the least-damped mode is not available as an alternative: at the readout rate $\kappa=0.958$ that mode exists only for $k<\kappa\pi/2=1.50$ (Section~\ref{sec:cut}), which carries about half of the profile weight ($0.51$ on the $36$ nodes used here, $0.48$ by direct integration); for the other half the pole has met the branch cut and there is no isolated kept eigenvalue. So, of the rulers tested, the law is a law in the operator's own link rapidity, and not in the displacement of the linear kept state: the difference between the two collision models is not carried by the linear kept state at the readout epoch \grade{measured; Appendix~B.17}.

\paragraph{Reading.} The anisotropic stress carries the threshold. Damping the shear faster than the flux keeps the stress from building and makes the collapse more fluid-like; damping it slower lets the stress the free-streaming effect is made of survive. The threshold reads the direction of that mismatch, in the rapidity of the collision operator's flux--shear link, and not how much energy sits on the link. That is the place Section~\ref{sec:witness} found the memory: in the angular structure a fluid closure drops. The kept-sector memory form, which counts only what the density keeps, cannot see it, and neither can the linear kept state at the first maximum \grade{H}.

\begin{figure}[t]\centering\includegraphics[width=\textwidth]{figs/fig7_link_v5.png}
\caption{The flux--shear link. Left: the surviving share $G$ against the link's reflection coefficient $\Gamma_{12}$, eleven thresholds from McCormick's code with the shear and higher moments relaxing at $r$ times the flux rate. Right: the offset of $G$ against the link rapidity $\tfrac12\ln r$, with the shear-only pairs (squares); the shaded band is the admissible window $\tfrac34\le r\le\tfrac32$. Lines: rate one (dashed) and the first-level inner-pair rate $1.11$.}\label{fig:link}\end{figure}

\paragraph{An independent check of a 3 November review.} An outside review of the companion paper reported two further thresholds on McCormick's code: a step in which collisions stop for good at $\epsilon_b=1$, at $\delta_c=0.513512$, and $n=3$, $\epsilon_{\rm dec}=1.20$, at $0.516797$. Here $0.5132$ and $0.5138$ bracket the first, and $0.5165$ and $0.5171$ the second, each dispersing below and forming a black hole above \grade{measured}.

\input{sec_new_gramian}

\input{sec_new_door}

\input{sec_new_mcc}

\section{Tests for McCormick's code}

\begin{enumerate}
\item \textbf{The annihilation control.} During e$^\pm$ annihilation the plasma sound speed falls to $c_s=0.519$ at $T=0.131$\,MeV, while the neutrinos' root-mean-square speed stays $1/\sqrt3$. The contest between front and sound, $f=(1-c_s)/(1-1/\sqrt3)$, reaches $1.135$ at $1.6\times10^7\,M_\odot$ ($c_s=0.521$ there). If the softer plasma enhanced the free-streaming effect, the shift would reach $7.5\%$, outside McCormick's envelope of $5.3$--$6.2\%$, so that reading already fails. If it reduces it, the shift stays at $5.5$--$6.3\%$ \grade{computed; reading}. At $1.6\times10^7\,M_\odot$ the two readings differ from the unmodified $6.31\%$ by $+13.5\%$ and $-11.9\%$. The run McCormick lists as missing, the same transport at fixed $R_\nu$ with $p=e/3$ against the full equation of state at $3\times10^6$--$1.6\times10^7\,M_\odot$, would confirm the sign.
\item \textbf{Which mean of the rates.} McCormick calibrates his single rate to an energy-group calculation. The arithmetic mean of the group rates reproduces his calibrated decoupling temperatures: $1.62$\,MeV against $1.6$ (e$^\pm$ only), $1.02$\,MeV against $0.93$--$1.0$ (with $\nu\nu$), and a mass step of $2.53$ against $2.5$--$3$. The harmonic and geometric means fail; the mean of squares also passes, so the class of means is identified, not one mean~\cite{WB} \grade{computed; reading}.
\end{enumerate}

\section{Relation to other work}

The closest published analogue is bulk viscosity, which raises the threshold roughly in proportion to its strength in full numerical relativity~\cite{Yuwen}. Free streaming is a second dissipative channel with the same sign. Papanikolaou fixes thresholds at horizon crossing for a time-dependent equation of state, with $\delta_c$ following $c_s^2$~\cite{Pap}. Escriv\`a ties them to the initial profile and the curvature of the compaction at its peak~\cite{Esc}. Test A is consistent with both: the neutrino clock is set early.

The 2026 thermal-history mass functions through neutrino decoupling and e$^\pm$ annihilation treat neutrinos as part of the fluid~\cite{Gonin}, so the step McCormick found is absent from them. For supermassive seeds, the $\mu$-distortion bound leaves only heavy-tailed statistics~\cite{Dave}, for which McCormick finds that the raised threshold costs a factor $1.2$--$8$ in abundance in his examples~\cite{McC}.

The nearest analogue to the result of Section~\ref{sec:link} is viscous damping. Bulk viscosity enters the collapse through the isotropic pressure~\cite{Yuwen}; the shear relaxation rate of a free-streaming component enters through the anisotropic stress, and no published threshold study varies it. In relaxation-time kinetic theory the kept sector's response is a continued fraction whose hydrodynamic pole meets a branch cut as the collision rate falls~\cite{Rom,Hu}; Section~\ref{sec:link} finds that at the readout scale the kept sector has no isolated eigenvalue for about half of the profile, so its eigenprojector cannot serve as a ruler there; of the rulers tested, the threshold is a law in the operator's link rapidity and not in the displacement of a kept state.

A search of the literature to October 2026 found no other work on thresholds with free-streaming components, on gradual kinetic decoupling in numerical relativity, on how the threshold depends on when a collision rate switches off, or on how it depends on the ratio of the shear and flux relaxation rates. The runs of Sections~\ref{sec:runs} and~\ref{sec:link} are, as far as that search shows, the first such runs.

McCormick's revised calculation~\cite{McC2} replaces the single rate by energy groups with plasma and self-scattering rates; Section~\ref{sec:mcc} shows that his collision term is exactly the split link of Section~\ref{sec:link}, so his model sits on the axis this paper varies, at $r\approx2.4$ for $\nu_e$ and $7.7$ for $\nu_{\mu,\tau}$. Froustey, Pitrou and Volpe solve neutrino decoupling in momentum space with full collision integrals and flavour mixing~\cite{FPV}; there the angular structure is that of a homogeneous background, while the hidden moments here are spatial. Seibold and Frank close moment models of radiative transfer by optimal prediction, with a memory term from the eliminated moments~\cite{SF}; Section~\ref{sec:gramian} finds that a short-memory correction of that type beats $P_2$ in every history tested, and that freezing the exact memory kernel to one history costs up to $0.13$. The readout Gramian is the observability Gramian of linear control theory; what is new here is its use as the weight on the predictive quotient of VOTKP~\cite{VOTKP} and Dirac Time~\cite{DT}, and the test of its leading direction in the full collapse.

\section{Summary}

The neutrino collision operator is a cut whose whole non-normality sits on one link, density to flux, with weight equal to the sound speed; its kept-sector mode lives only above Romatschke's onset $\kappa=2k/\pi$, where its non-orthogonality diverges. McCormick's code, run on cases outside his paper, shows that the collapse time alone does not set the neutrino clock, and that the decoupling history changes the threshold, not only its end point. The kept-sector memory, read mostly at the first compaction maximum and fitted once, predicted three histories it had not seen to within $0.071$, $0.014$ and $0.037$, while a history-blind rule failed. The steep-history thresholds do not move with the rate cap or the grid. The memory sits in the angular structure a fluid closure drops: set at or after decoupling, it returns into density, flux and stress; set deep in the coupled era, it is erased. And the anisotropic stress carries the threshold: with the shear relaxing at $r$ times the flux rate, the offset of the surviving share moves close to a straight line in the link rapidity $\tfrac12\ln r$, at a rate of $1.10\pm0.01$ inside the admissible window, and at $1.19$--$1.20$ when the shear alone is damped. The direction of the flux--shear mismatch, not its connection energy, is what the collapse reads; of the rulers tested, the law lives in the collision operator's own link rapidity, not in the displacement of the linear kept state; and the rate is a measured number at finite resolution, $1.10\pm0.01$, the $\pm0.01$ an error budget rather than a converged value or a statistical interval, $1.102$ and $1.101$ at the two finest of three grid levels, with no closed form attached. The readout Gramian of the hidden moments has no kernel and ranks what must be kept, and collision histories that agree in integrated opacity and integrated connection energy still differ in the memory. In the full collapse, hidden angular structure planted at decoupling, invisible to density, flux, stress and Einstein's equations at the restart, moves the black-hole threshold by up to $5.5\times10^{-3}$, with the sign set by its parity and the size ranked by the Gramian; planted in the coupled era it moves it by less than one per cent of that. McCormick's physical collision term is the split link itself, and self-scattering acts on the threshold mostly through timing, with the flux--shear split carrying $+0.0003$. What the neutrino state must keep is not its density, flux and stress alone, but the hidden angular structure the future reads, and the order of the history that made it.

\begin{thebibliography}{99}
\bibitem{McC} S.~McCormick, \emph{Free-streaming neutrinos raise the threshold for primordial black hole formation}, preprint (2026). Simulation code, threshold data and scripts: Zenodo, doi:10.5281/zenodo.23808375.
\bibitem{McCcode} S.~McCormick, PBH\_neutrino\_simulation (code and threshold data), github.com/quasilocal/PBH\_neutrino\_simulation (2026); archived as doi:10.5281/zenodo.23808375.
\bibitem{McC2} S.~McCormick, \emph{Free-streaming neutrinos raise the threshold for primordial black hole formation}, version 1.1, with multigroup decoupling (2026). Code: github.com/quasilocal/PBH\_neutrino\_simulation, commit eb7c4de; Zenodo, doi:10.5281/zenodo.24129915.
\bibitem{FPV} J.~Froustey, C.~Pitrou and M.~C.~Volpe, \emph{Neutrino decoupling including flavour oscillations and primordial nucleosynthesis}, JCAP 12 (2020) 015, arXiv:2008.01074.
\bibitem{SF} B.~Seibold and M.~Frank, \emph{Optimal prediction for moment models: crescendo diffusion and reordered equations}, Contin.\ Mech.\ Thermodyn.\ 21 (2009) 511, arXiv:0902.0076.
\bibitem{VOTKP} J.~Beasley, \emph{VOTKP: The Vector of Traits Kept under Peeling}, v2.9 (2026), Zenodo, doi:10.5281/zenodo.23340757.
\bibitem{DT} J.~Beasley, \emph{The Dirac Time of the Gigantefermion: Memory First, Projector Geometry, Retained History, and the Cost of an Arrow}, v3.0.1 (2026), Zenodo, doi:10.5281/zenodo.22978631.
\bibitem{CE} J.~Beasley, \emph{Non-Normality Is Connection Energy} (2026), Zenodo, doi:10.5281/zenodo.22804072.
\bibitem{WB} J.~Beasley, \emph{The Weight Between: What the present must keep, and what it costs to move it}, v1.1 (2026).
\bibitem{OSL} J.~Beasley, \emph{The Offset Speed Limit: Rate is weight}, v1.1 (2026).
\bibitem{CSL} J.~Beasley, \emph{Colour Selects the Leader}, v1.1 (2026).
\bibitem{BH} J.~Beasley, \emph{Beyond the Event Horizon: What Survives at a Rigidity Wall}, edition 10.1 (2026).
\bibitem{Rom} P.~Romatschke, \emph{Retarded correlators in kinetic theory: branch cuts, poles and hydrodynamic onset transitions}, Eur.\ Phys.\ J.\ C 76 (2016) 352, arXiv:1512.02641.
\bibitem{Hu} J.~Hu, \emph{Relaxation time approximation revisited and non-analytical structure in retarded correlators}, SciPost Phys.\ 20 (2026) 020, arXiv:2409.05131.
\bibitem{Yuwen} Yuwen, Joana, Wang and Cai, \emph{PBH formation with bulk viscosity}, arXiv:2606.26532 (2026).
\bibitem{Pap} T.~Papanikolaou, \emph{Toward the primordial black hole formation threshold in a time-dependent equation-of-state background}, arXiv:2205.07748.
\bibitem{Esc} A.~Escriv\`a, arXiv:2504.05814.
\bibitem{Gonin} Gonin, Froustey, Escriv\`a, Magaraggia, K\"uhnel and Hasinger, \emph{Primordial asymmetries, primordial equation of state and primordial black holes}, arXiv:2608.22240 (2026).
\bibitem{Dave} Dave, Yan, Ilyas and Cai, \emph{Evading the CMB $\mu$-distortion bound on supermassive PBH seeds with non-Gaussian tails}, arXiv:2607.03138 (2026).
\bibitem{MB} C.-P.~Ma and E.~Bertschinger, Astrophys.\ J.\ 455, 7 (1995).
\bibitem{IQGT} J.~Beasley, \emph{Intrinsic Quantum Geometric Tensor on the Grassmannian and Spectral Lower Bounds on Dissipation}, v7.2, Zenodo, doi:10.5281/zenodo.23351266 (2026).
\bibitem{Atlas} J.~Beasley, \emph{The Relational Atlas} and \emph{The Quiet Law} (drafts, 2026), in the VOTKP v2.9 supporting records, Zenodo, doi:10.5281/zenodo.23340757.
\bibitem{SMball} J.~Beasley, \emph{The Standard Model in the ball, and what happens at the walls}, working notes v3, 11 November 2026.
\bibitem{Alik} I.~Alikhanov, \emph{Another relation among the neutrino mass-squared differences?}, arXiv:2601.18781 (2026).
\bibitem{CCC} R.~Penrose, \emph{Cycles of Time} (Bodley Head, 2010); V.~G.~Gurzadyan and R.~Penrose, Eur.\ Phys.\ J.\ Plus 128 (2013) 22.
\bibitem{ST} P.~J.~Steinhardt and N.~Turok, Science 296 (2002) 1436; Phys.\ Rev.\ D 65 (2002) 126003.
\end{thebibliography}

\appendix
\section{Proofs}
\subsection*{A.1 Non-normality is connection energy, for the neutrino chain}
Let $L=A+D$ with $A^\ast=-A$ and $D^\ast=D$. Then $L^\ast=D-A$, and
\[
[L,L^\ast]=(A+D)(D-A)-(D-A)(A+D)=2(AD-DA)=2[A,D].
\]
For diagonal $D=\mathrm{diag}(D_\ell)$, $[A,D]_{ij}=A_{ij}(D_j-D_i)$. So
\[
\tfrac18\lVert[L,L^\ast]\rVert_F^2=\tfrac12\sum_{i\neq j}|A_{ij}|^2(D_i-D_j)^2 .
\]
Streaming couples only $\ell$ and $\ell\pm1$, with $|A_{\ell,\ell+1}|=kc_\ell$. Each link appears twice in the ordered sum, which gives $\sum_\ell k^2c_\ell^2(D_{\ell+1}-D_\ell)^2$. With $D=(0,-\kappa_1,-\kappa_2,-\kappa_2,\dots)$, $c_0^2=1/3$ and $c_1^2=4/15$, this is $k^2\kappa_1^2/3+\tfrac{4k^2}{15}(\kappa_2-\kappa_1)^2$. With $\kappa_1=\kappa_2=\kappa$ it is $\kappa^2k^2/3$. $\square$
\subsection*{A.2 The weight across the cut}
$PH(I-P)$ with $H=k\mu$ has the single entry $\langle0|k\mu|1\rangle=kc_0$, so $E_B=k^2/3$. The variance of $k\mu$ in the isotropic state is $k^2\langle\mu^2\rangle=k^2/3$, the same number. $\square$
\subsection*{A.3 The kept-sector propagator}
Write $u=(s-L)^{-1}|0\rangle$, so $(s+\kappa+ik\mu)u=|0\rangle+\kappa Pu$. Projecting onto $|0\rangle$ with $\langle0|(s+\kappa+ik\mu)^{-1}|0\rangle=\tfrac12\int_{-1}^{1}d\mu/(s+\kappa+ik\mu)=k^{-1}\arctan[k/(s+\kappa)]=g$ gives $u_0=g(1+\kappa u_0)$, so $u_0=g/(1-\kappa g)$. $\square$
\subsection*{A.4 The truncated exceptional point}
On $\mathrm{span}\{|0\rangle,|1\rangle\}$, $L=\begin{pmatrix}0&-ik/\sqrt3\\-ik/\sqrt3&-\kappa\end{pmatrix}$, with characteristic polynomial $\lambda^2+\kappa\lambda+k^2/3$. Its discriminant $\kappa^2-4k^2/3$ vanishes at $\kappa=2k/\sqrt3$, where the matrix is not diagonalisable. $\square$
\subsection*{A.3a The onset of the kept-sector pole (Romatschke's onset~\cite[Eq.~(37)]{Rom}, with its non-orthogonality)}
A pole of $u_0=g/(1-\kappa g)$ needs $\kappa g=1$, that is $\arctan[k/(s+\kappa)]=\phi$ with $\phi=k/\kappa$. For real $s>-\kappa$ the principal arctangent takes values in $(0,\pi/2)$, so a solution exists exactly when $\phi<\pi/2$, and then $s+\kappa=k\cot\phi$. Because the principal arctangent is one-to-one, there is no other solution on the physical sheet. As $\phi\to\pi/2$, $s\to-\kappa$, the centre of the cut of $g$, which joins the branch points $-\kappa\pm ik$. The eigenvector is $u(\mu)=1/(s+\kappa+ik\mu)$. The operator is complex symmetric in the Legendre basis, so its left eigenvector is $u$ itself under the bilinear pairing, and $K=\big(\tfrac12\!\int|u|^2\big)^2/\big|\tfrac12\!\int u^2\big|^2$. With $a=s+\kappa=k\cot\phi$,
\[
\tfrac12\int_{-1}^{1}\frac{d\mu}{a^2+k^2\mu^2}=\frac{\arctan(k/a)}{ak}=\frac{\phi\tan\phi}{k^2},\qquad
\tfrac12\int_{-1}^{1}\frac{d\mu}{(a+ik\mu)^2}=\frac{1}{a^2+k^2}=\frac{\sin^2\phi}{k^2},
\]
so $K=\phi^2\tan^2\phi/\sin^4\phi=\phi^2/(\sin^2\phi\cos^2\phi)=4\phi^2/\sin^22\phi$. $\square$

\subsection*{A.5 Exact elimination of the hidden moments}
Write $\partial_\tau v=-(ikM+\sigma(\tau)Q)v$ with $Q=I-P$, and split $v=(a,h)$ along $P$ and $Q$. Then $\dot a=-ikc_0h_1$ and $\dot h=-(ikM_{QQ}+\sigma)h-ikc_0\,a\,e_1$. On the hidden block $\sigma(\tau)I$ commutes with $M_{QQ}$, so the generators at different times commute and the propagator is the product $U_Q(t,r)=e^{-\int_r^t\sigma}\,e^{-ik(t-r)M_{QQ}}$, with $\int_r^t\sigma=\Sigma(t)-\Sigma(r)$, $\Sigma(\tau)=-\epsilon_{\rm dec}^n/(n\tau^n)$. Variation of constants gives $h(t)=U_Q(t,s)h(s)-ikc_0\int_s^tU_Q(t,r)e_1a(r)\,dr$; substituting into $\dot a$ gives
\[
\dot a(t)=-ikc_0[U_Q(t,s)h(s)]_1-k^2c_0^2\int_s^t[U_Q(t,r)]_{11}a(r)\,dr .
\]
For two states with equal $a(s)$ the difference $\Delta a$ obeys the same equation with source $S[\Delta h(s)]$ and $\Delta a(s)=0$; if $\Delta h(s)=0$ it vanishes identically. The time-independent case $e^{t(X-\sigma I)}=e^{-\sigma t}e^{tX}$ is checked in Lean (Appendix~C). $\square$
\subsection*{A.6 The speed limit on the kept population}
With $D=-\sigma Q\le0$, $\frac{d}{d\tau}\lVert v\rVert^2=-2\sigma\lVert Qv\rVert^2\le0$, so $\lVert v\rVert\le1$ for a unit initial state. Since $\langle0|D=0$ and $\langle0|M=c_0\langle1|$, $\dot v_0=-ikc_0v_1$, and
\[
\dot A=2\,\mathrm{Re}\big(\bar v_0\,\dot v_0\big)=2\,\mathrm{Re}\big(\bar v_0(-ikc_0)v_1\big),\qquad |\dot A|\le2kc_0|v_0||v_1|.
\]
With $|v_1|^2\le\lVert v\rVert^2-|v_0|^2\le1-A$ and $k^2c_0^2=E_B$, $\dot A^2\le4E_BA(1-A)$. The core inequality $|2\,\mathrm{Re}(\bar a(-ih)b)|\le2|h||a||b|$ and the corollary under $|a|^2+|b|^2\le1$ are checked in Lean (Appendix~C). $\square$
\subsection*{A.7 Energy stationarity is not response stationarity}
Keep three moments, $\ell=0,1,2$, with streaming weights $c_0^2=1/3$, $c_1^2=4/15$, and $L=-ikM-D$ with $D=\mathrm{diag}(0,\kappa_1,\kappa_2)$, so that $s-L=s+ikM+D$ (the damping enters $L$ with a minus sign, as in Section~\ref{sec:cut}). The kept response $R_{00}=\langle0|(s-L)^{-1}|0\rangle$ is the continued fraction
\[
R_{00}=\cfrac{1}{s+\cfrac{k^2c_0^2}{s+\kappa_1+\cfrac{k^2c_1^2}{s+\kappa_2}}} .
\]
At $k=s=\kappa_1=1$ this is $R_{00}(\kappa_2)=\big(1+\tfrac13/(2+\tfrac{4/15}{1+\kappa_2})\big)^{-1}$, so $R_{00}(1)=32/37$ and $R_{00}'(1)=-5/1369$. The energy $E=k^2\kappa_1^2/3+\tfrac{4k^2}{15}(\kappa_2-\kappa_1)^2$ has $\partial E/\partial\kappa_2=\tfrac{8k^2}{15}(\kappa_2-\kappa_1)=0$ there. Checked by direct inversion of the $3\times3$ matrix; checked in Lean (Appendix~C). $\square$

\subsection*{A.8 The link coordinate is odd}
For $r>0$, $\Gamma(1/r)=(1/r-1)/(1/r+1)=(1-r)/(1+r)=-\Gamma(r)$, and $\tfrac12\ln(1/r)=-\tfrac12\ln r$. Hence $\Gamma(r)=\tanh(\tfrac12\ln r)$, since $\tanh u=(e^{2u}-1)/(e^{2u}+1)$ with $e^{2u}=r$. $\square$

\subsection*{A.9 What cancels from the odd part, and what does not}
Write $\tau_G(r)=\ln[(1-G(r))/G(r)]$ and let $\tau_1=\tau_G(1)$. The odd part of the pair $r,1/r$ is $[(\tau_G(r)-\tau_1)-(\tau_G(1/r)-\tau_1)]/\ln r=[\tau_G(r)-\tau_G(1/r)]/\ln r$: the anchor $\tau_1$ cancels identically. With $G=(\delta-\delta_f)/(\delta_i-\delta_f)$, $\tau_G=\ln(\delta_i-\delta)-\ln(\delta-\delta_f)$, so
\[
\frac{\partial}{\partial\delta_f}\big[\tau_G(r)-\tau_G(1/r)\big]=\frac1{\delta_r-\delta_f}-\frac1{\delta_{1/r}-\delta_f},\qquad
\frac{\partial}{\partial\delta_i}\big[\tau_G(r)-\tau_G(1/r)\big]=-\frac1{\delta_i-\delta_r}+\frac1{\delta_i-\delta_{1/r}},
\]
which vanish only if $\delta_r=\delta_{1/r}$. The reference thresholds therefore enter the odd rate, with coefficients $\partial(\text{rate})/\partial\delta_f\approx+50$ and $\partial(\text{rate})/\partial\delta_i\approx-28$ per unit for the inner pairs, which is the $\pm0.004$ common-mode term of Section~\ref{sec:link}. $\square$

\subsection*{A.10 The streaming weights}
$c_\ell^2=(\ell+1)^2/((2\ell+1)(2\ell+3))$ gives $c_0^2=1/3$, $c_1^2=4/15$, $c_2^2=9/35$, so $c_0^2/c_1^2=5/4$ and $c_1^2/c_2^2=28/27$: the adjacent weights differ by $1.8\%$ from the second link on, which is why a rule additive in link-weight ratios nearly cancels when two adjacent links jump by opposite rapidities (Section~\ref{sec:link}). Checked in Lean. $\square$

\subsection*{A.11 Every verdict as a whole-number inequality}
Let $O(\delta)=(\delta_i-\delta)/(\delta-\delta_f)$, the odds $(1-G)/G$, so $\tau_G=\ln O$, and let $Q=O(\delta_r)/O(\delta_{1/r})$. The odd rate is $\ln Q/\ln r$, and for $r=u/v$ and a rational bound $p/q$,
\[
\frac{\ln Q}{\ln r}\ge\frac pq\iff Q^q\ge r^p\iff \big[(\delta_i-\delta_r)(\delta_{1/r}-\delta_f)\big]^q v^p\ge\big[(\delta_r-\delta_f)(\delta_i-\delta_{1/r})\big]^q u^p ,
\]
with every threshold written as an integer on a common scale. $O$ decreases in $\delta$, so over a pair of brackets $Q$ is smallest at the upper corner of the $r$ bracket with the lower corner of the $1/r$ bracket, and largest at the opposite pair of corners; checking both corners bounds the rate over the whole bracket pair. Each such inequality is decided exactly by Lean's kernel (Appendix~C). $\square$

\subsection*{A.12 The discrete dipole and shear projections}
On the midpoint bins $\mu_k$ with weight $\Delta\mu$, write $\langle f,g\rangle=\sum_kf_kg_k\Delta\mu$, $m_2=\langle\mu,\mu\rangle$, and $P_2=\mu^2-m_2/2$. Then $\langle1,\mu\rangle=0$ by the symmetry of the bins, $\langle1,P_2\rangle=m_2-\tfrac{m_2}{2}\langle1,1\rangle=0$ because $\langle1,1\rangle=2$, and $\langle\mu,P_2\rangle=0$ by symmetry. So $d=\langle\mu,I\rangle\mu/m_2$ and $s=\langle P_2,I\rangle P_2/\langle P_2,P_2\rangle$ are the orthogonal projections of the intensity onto the discrete $\ell=1$ and $\ell=2$ directions, orthogonal to each other and to the isotropic part. The collision terms of Section~\ref{sec:link}, $\kappa_2(\varepsilon/4\pi-I)+(\kappa_2-\kappa_1)d$ and $\kappa_1(\varepsilon/4\pi-I)+(\kappa_1-\kappa_2)s$, therefore conserve the bin-summed energy exactly, leave the dipole relaxing at $\kappa_1$ in both, and relax the $\ell=2$ part at $\kappa_2$ in both (together with every $\ell\ge3$ in the first, alone in the second); $m_2=2/3-\Delta\mu^2/6$ differs from the continuum $2/3$ by $0.1\%$ at $N_\mu=32$. Checked numerically to $10^{-16}$ (Appendix~B.15). $\square$

\input{app_new_proofs}

\section{Evidence: scripts, run logs and outputs}
\subsection*{B.1 Identity checks}
\lstinputlisting{code/main/identities.py}
\lstinputlisting{code/main/identities_out.txt}
\subsection*{B.2 The one change to McCormick's code}
In a copy of \texttt{run\_nuc.py}, the collision rate line became\\
\texttt{kap = (1/(2*t))*min((eps\_dec/(2*np.sqrt(t*t0)))**nexp, kapcap)}\\
with \texttt{nexp} (default 3) and \texttt{kapcap} $=5\times10^5$ as inputs. Nothing else was changed.
\lstinputlisting{code/bisect_beyond.py}
\lstinputlisting{code/trisect_beyond.py}
\subsection*{B.3 Predictions written before the runs}
\lstinputlisting{code/PREREGISTERED.txt}
\subsection*{B.4 Run logs}
\lstinputlisting{results/log_repro.txt}
\lstinputlisting{results/log_A4.txt}
\lstinputlisting{results/log_A8.txt}
\lstinputlisting{results/log_A8b.txt}
\lstinputlisting{results/log_B2.txt}
\lstinputlisting{results/log_B3.txt}
\subsection*{B.5 Kept-sector memory curves}
\lstinputlisting{code/main/memory_curve.py}
\lstinputlisting{code/main/memory_curve_out.txt}
\subsection*{B.6 The kept-sector readings against McCormick's curve and the history prediction}
\lstinputlisting{code/wb_checks/wb_steve_03_paper_objects.py}
\lstinputlisting{code/wb_checks/wb_steve_03_out.txt}
\lstinputlisting{code/wb_checks/wb_steve_07_history.py}
\lstinputlisting{code/wb_checks/wb_steve_07_out.txt}
\lstinputlisting{results/provenance_wb07.txt}
\subsection*{B.7 Shear and flux links, and which mean of the rates}
\lstinputlisting{code/wb_checks/wb_steve_01_algebra.py}
\lstinputlisting{code/wb_checks/wb_steve_01_out.txt}
\lstinputlisting{code/wb_checks/wb_steve_04_channels.py}
\lstinputlisting{code/wb_checks/wb_steve_04_out.txt}
\subsection*{B.8 The full operator and the truncated wall}
\lstinputlisting{code/osl_checks/osl_b_ep_full_operator.py}
\lstinputlisting{code/osl_checks/osl_b_out.txt}
\subsection*{B.9 The annihilation contest}
\lstinputlisting{code/osl_checks/osl_c_speed_contest_eos.py}
\lstinputlisting{code/osl_checks/osl_c_out.txt}
\subsection*{B.10 Linear first maximum under gradual decoupling}
\lstinputlisting{code/linear_solver/lin.py}
\lstinputlisting{code/linear_solver/grad.py}
\lstinputlisting{code/linear_solver/grad.out}
\subsection*{B.11 Rate-ratio and wall readings}
\lstinputlisting{code/fixed.py}
\lstinputlisting{code/fixed_out.txt}
\lstinputlisting{code/main/lone_pass_check.py}
\lstinputlisting{code/main/lone_pass_check_out.txt}
\subsection*{B.12 Diffusive escape}
\lstinputlisting{code/csl_checks/k6_diffusion_leader.py}
\lstinputlisting{code/csl_checks/k6_diffusion_leader_out.txt}
\subsection*{B.13 Tests C and D and the controls: predictions, search script and run logs}
\lstinputlisting{code/v3_runs/PREREGISTERED_NEXT.txt}
\lstinputlisting{code/v3_runs/trisect3.py}
\lstinputlisting{results/v3/log_Bcap_part1.txt}
\lstinputlisting{results/v3/log_Bcap2.txt}
\lstinputlisting{results/v3/log_Bres.txt}
\lstinputlisting{results/v3/log_Bres2.txt}
\lstinputlisting{results/v3/log_C.txt}
\lstinputlisting{results/v3/log_D.txt}
\subsection*{B.14 The matched-pair witness}
\lstinputlisting{code/witness/PREREGISTERED_WITNESS.txt}
\lstinputlisting{code/witness/witness_run.py}
\lstinputlisting{code/witness/witness_kernel.py}
\lstinputlisting{code/witness/witness_conv.py}
\lstinputlisting{code/witness/witness_kernel.json}
\lstinputlisting{code/witness/prod.log}
\lstinputlisting{code/witness/convLM.log}
\lstinputlisting{code/witness/convTol.log}
\subsection*{B.15 The flux--shear link and the neutrino fraction: predictions, changed solvers, drivers, run logs}
Every prediction and kill line of Section~\ref{sec:link} was written before the run it judges; the full pre-registration file, with each result appended under its test and including the intermediate readings that the final runs superseded, is \texttt{code/v4\_runs/PREREGISTERED\_V4.txt} in the archive and on GitHub. The final runs are the ones listed here. The changed solvers: McCormick's \texttt{ms\_nuc.py} (MIT licence, copyright S.~McCormick; the licence file ships with the archive) with the collision term split, and the shear-only variant; each differs from McCormick's file in its header line, the signature of \texttt{rhs}, the lines that compute the discrete projection and $\kappa_2$, and the collision term. The runners label a breakdown after $2M/R>0.95$ separately from a detected horizon and log the minimum intensity; the driver moves a bracket only on a detected horizon or a dispersal.
\lstinputlisting{code/v4_runs/ms_nuc_split.py}
\lstinputlisting{code/v4_runs/ms_nuc_shear.py}
\lstinputlisting{code/v4_runs/run_split_v2.py}
\lstinputlisting{code/v4_runs/trisect_v2.py}
\lstinputlisting{code/v4_runs/trisect_clean.py}
\lstinputlisting{code/v4_runs/trisect_split.py}
\lstinputlisting{code/v4_runs/trisect_fnu.py}
\lstinputlisting{code/v4_runs/memory_split.py}
\lstinputlisting{code/v4_runs/log_repro.txt}
\lstinputlisting{code/v4_runs/log_F_regress.txt}
\lstinputlisting{code/v4_runs/log_L_regress.txt}
\lstinputlisting{code/v4_runs/log_E.txt}
\lstinputlisting{code/v4_runs/log_F025.txt}
\lstinputlisting{code/v4_runs/log_F4.txt}
\lstinputlisting{code/v4_runs/log_F2a.txt}
\lstinputlisting{code/v4_runs/log_F2b.txt}
\lstinputlisting{code/v4_runs/log_H1.txt}
\lstinputlisting{code/v4_runs/log_I.txt}
\lstinputlisting{code/v4_runs/log_J2.txt}
\lstinputlisting{code/v4_runs/log_J05.txt}
\lstinputlisting{code/v4_runs/log_K125_clean.txt}
\lstinputlisting{code/v4_runs/log_K08_clean.txt}
\lstinputlisting{code/v4_runs/log_L4.txt}
\lstinputlisting{code/v4_runs/log_L025.txt}
\lstinputlisting{code/v4_runs/log_M4.txt}
\lstinputlisting{code/v4_runs/log_M025.txt}
\lstinputlisting{code/v4_runs/log_A1_shear_r12.txt}
\lstinputlisting{code/v4_runs/log_A1_shear_r083.txt}
\lstinputlisting{code/v4_runs/log_A1_all_r12.txt}
\lstinputlisting{code/v4_runs/log_A1_all_r083.txt}
The endpoint audit: every decisive bracket edge of every threshold above, rerun with the labelled runners.
\lstinputlisting{code/v4_runs/log_B1_edges.txt}
\subsection*{B.16 The check of the 3 November review}
\lstinputlisting{code/v4_runs/run_step.py}
\lstinputlisting{code/v4_runs/queue_G.sh}
\lstinputlisting{code/v4_runs/log_G.txt}
\subsection*{B.17 The coordinate of the law: the intrinsic-ruler check}
\lstinputlisting{code/v4_runs/ruler/PREREGISTERED_RULER.txt}
\lstinputlisting{code/v4_runs/ruler/ruler_check.py}
\lstinputlisting{code/v4_runs/ruler/diag_modes.py}
\lstinputlisting{code/v4_runs/ruler/ruler_state.py}
\lstinputlisting{code/v4_runs/ruler/out_state.txt}
\subsection*{B.18 Figures}
\lstinputlisting{code/main/make_figures.py}
\lstinputlisting{code/main/make_figures_v3.py}
\lstinputlisting{code/main/make_fig6_door.py}
\lstinputlisting{code/v4_runs/make_fig7_link_v5.py}
\lstinputlisting{code/v4_runs/make_fig7_link_v5_out.txt}
\input{app_new_evidence}

\section{Lean statements}
\sloppy
The statements below are proved in Lean~4 with Mathlib (v4.34.1) in the repository \url{https://github.com/dicipler-pixel/neutrinos-leave-lean}. A full build of the Lean sources at commit \texttt{9cd71fa} (13 November 2026, 8934 build jobs) completes with no \texttt{sorry}; the audit script (commit \texttt{9292b2c}, which changes only its configuration) then finds that all 87 theorems of the library use only the standard axioms (\texttt{propext}, \texttt{Quot.sound}, \texttt{Classical.choice}) and that all eleven false controls are rejected for a mathematical reason. The log, with the toolchain and Mathlib pins, ships with the archive. An earlier build at commit \texttt{c372d3f} (4 November 2026) replayed every proof in an independent kernel checker. Lean checks the finite moment model and the verdict arithmetic. It does not check McCormick's code, the collapse runs, the classification of runs near threshold, the bracket edges (inputs copied from the run logs) or the numerical witness.
\begin{enumerate}
\item \texttt{crossLink}: $|2\,\mathrm{Re}(\bar a(-ih)b)|\le2|h||a||b|$; \texttt{crossLink\_sq}: with $A=|a|^2$, $|a|^2+|b|^2\le1$, the square is at most $4|h|^2A(1-A)$ (Appendix A.6).
\item \texttt{twoMomentConnectionEnergy}: $\tfrac18\lVert[L,L^\ast]\rVert_F^2=\kappa^2k^2/3$ for the two-moment $L$.
\item \texttt{splitRateEnergy}: $\tfrac18\lVert[L,L^\ast]\rVert_F^2=k^2\kappa_1^2/3+\tfrac{4k^2}{15}(\kappa_2-\kappa_1)^2$ for the three-moment $L$.
\item \texttt{shearCounterexample}: $R_{00}(1)=32/37$, $R_{00}'(1)=-5/1369$, $\partial_{\kappa_2}E=0$ at $\kappa_2=\kappa_1$; with the complex $3\times3$ resolvent entry equal to the closed form $2(15\kappa_2+17)/(35\kappa_2+39)$ for $\kappa_2>-1$ (Appendix A.7).
\item \texttt{scalarCollisionPropagator}: $e^{t(X-\sigma I)}=e^{-\sigma t}e^{tX}$, with generators $X-\sigma(t)I$ commuting pairwise (Appendix A.5).
\item \texttt{testB/C/D\_verdict}: $G_B=970/3059$, $G_C=1983/3059$, $G_D=472/3059$ against the pre-registered kill lines.
\item \texttt{gam\_inv}, \texttt{odd\_part\_free\_of\_anchor}: the link coordinate $(r-1)/(r+1)$ is odd under $r\mapsto1/r$, and the anchor cancels from the odd part (Section~\ref{sec:link}).
\item \texttt{cSq\_zero}, \texttt{cSq\_one}, \texttt{cSq\_two}, \texttt{weight\_ratio\_01}, \texttt{weight\_ratio\_12}: $c_\ell^2=(\ell+1)^2/((2\ell+1)(2\ell+3))$ gives $1/3$, $4/15$, $9/35$, and the ratios $5/4$ and $28/27$.
\item \texttt{testK\_clean\_range}, \texttt{testK\_clean\_below\_sqrt}, \texttt{testK\_clean\_above\_trend}: at $r=1.25$ against $0.8$ the odd rate lies in $[1.112,1.114]$ over the whole bracket pair; \texttt{testK\_crash\_edge\_false\_kill}, \texttt{testK\_first\_sqrt\_window}, \texttt{testK\_first\_trend\_not\_killed}: what the same pair would have decided had a failed run been used as a bracket edge (it would have decided wrongly), which is why the search rule of Section~\ref{sec:link} never lets one.
\item \texttt{testJ\_straddles}: over its full bracket the $r=2$ odd rate straddles $1.108$ and stays at most $1.128$.
\item \texttt{selfScatterEnergy}, \texttt{selfScatterEnergy\_ab}, \texttt{selfScatterEnergy\_no\_nn} (new in version 6, module \texttt{Neutrinos.SelfScatter}): with $\kappa_1=\kappa_e$ and $\kappa_2=\kappa_e+\kappa_n$, $\tfrac18\lVert[L,L^\ast]\rVert_F^2=a^2\kappa_e^2+b^2\kappa_n^2$ for any streaming weights, $k^2\kappa_e^2/3+\tfrac{4k^2}{15}\kappa_n^2$ in the continuum weights, and $a^2\kappa_e^2$ without self-scattering (Appendix A.13). The discrete weights $c_0^2=\tfrac13-\Delta\mu^2/12$, $c_1^2=\tfrac4{15}(1-\Delta\mu^2)$ and the diagonalisation of McCormick's operator are not checked in Lean.
\item \texttt{testL\_rate\_ge\_six\_fifths}, \texttt{testL\_rate\_le}: the shear-only odd rate lies in $[1.2,1.21]$.
\end{enumerate}
Each verdict is a whole-number inequality: with $O(\delta)=(\delta^{\rm inst}-\delta)/(\delta-\delta^{\rm fluid})$ and $Q=O(\delta_r)/O(\delta_{1/r})$, the odd rate is at least $p/q$ exactly when $r^p\le Q^q$; $O$ decreases in $\delta$, so the two corners of the bracket pair that make $Q$ smallest and largest are the ones checked, and the kernel decides each inequality exactly.


\end{document}
